\exercisesection{赋范代数}

\begin{enumerate}
\def\labelenumi{\arabic{enumi})}
\tightlist
\item
  设 (A, e) 为一个复幺代数。
\end{enumerate}

\begin{enumerate}
\def\labelenumi{\alph{enumi})}
\tightlist
\item
  对于线性型 \(\lambda \in \mathrm{A}^*\)，当 \(\lambda(e) = 1\) 时，下列条件等价：
\end{enumerate}

\begin{enumerate}
\def\labelenumi{(\roman{enumi})}
\item
  对每个 \(x \in \operatorname{Ker}(\lambda)\)，有 \(x^2 \in \operatorname{Ker}(\lambda)\)；
\item
  对每个 \(x \in \mathbb{A}\)，有 \(\lambda(x^{2}) = \lambda(x)^{2}\)；
\item
  \(\lambda\) 的核是 A 的一个右理想；
\item
  线性型 \(\lambda\) 是从 A 到 C 的保单位元态射。
\end{enumerate}

（注意 \(\lambda(x - \lambda(x)e) = 0\) 对所有 \(x\in A\) 成立。）

\begin{enumerate}
\def\labelenumi{\alph{enumi})}
\setcounter{enumi}{1}
\item
  设 \(\lambda \in \mathrm{A}^*\) 为满足 \(\lambda(1) = 1\) 的线性型，且使得对每个 \(x \in \mathrm{A}\)，有 \(\lambda(x) \in \mathrm{Sp}_{\mathrm{A}}(x)\)。设 \(x \in \mathrm{Ker}(\lambda)\) 满足 \(\lambda(x^{2}) \neq 0\)，则 \(x\) 的谱在 \(\mathbf{C}\) 中无界。（取整数 \(n \geqslant 2\)，考虑多项式 \(p \in \mathbf{C}[\mathrm{X}]\)，使 \(p(z) = \lambda((ze - x)^{n})\) 对所有 \(z \in \mathbf{C}\) 成立；注意它的根都属于 \(x\) 的谱，并用 \(|\lambda(x^{2})|\) 给出它们的欧氏范数的下界。）
\item
  假设 A 是 Banach 代数。为使 A 上的线性型 \(\lambda\) 是从 A 到 C 的保单位元态射，充要条件是 \(\lambda(e)=1\)，且 \(\lambda(x)\in\mathrm{Sp}_{\mathrm{A}}(x)\) 对所有 \(x\in A\) 成立。
\end{enumerate}

\begin{enumerate}
\def\labelenumi{\arabic{enumi})}
\setcounter{enumi}{1}
\tightlist
\item
  设 \(n \geqslant 0\) 为整数。设 \(A_n\) 为由函数 \(f: [0,1] \to \mathbf{C}\)（它在 \([0,1]\) 上具有直到 \(n\) 阶的连续导数）所组成的交换复 Banach 代数，赋以范数
\end{enumerate}

\[
\| f \| = \sum_ {k = 0} ^ {n} \frac {1}{k !} \sup _ {0 \leqslant t \leqslant 1} | f ^ {(k)} (t) |
\]

（I，第 18 页，例 4）。证明 \(A_{n}\) 同构于 \(A_{m}\) 当且仅当 n = m。（不过，根据 I，第 144 页，例 1，这些 Banach 代数都有相同的特征标空间。）

\begin{enumerate}
\def\labelenumi{\arabic{enumi})}
\setcounter{enumi}{2}
\tightlist
\item
  设 G 为等于其导出子群的群（A，I，第 67 页）。对每个 \(g \in G\)，用 \(\ell(g)\) 表示满足下列条件的最小整数 \(k \geqslant 0\)：存在数对 \((a_{1}, b_{1}), \ldots, (a_{k}, b_{k})\)（取自 \(G \times G\)），使 \(g = [a_{1}, b_{1}] \cdots [a_{k}, b_{k}]\)。
\end{enumerate}

\begin{enumerate}
\def\labelenumi{\alph{enumi})}
\item
  映射 \(\ell\) 是 G 上的一个度量。
\item
  对每个 \(g \in \mathbf{G}\)，极限
\end{enumerate}

\[
\operatorname{lsc} (g) = \lim _ {n \to + \infty} \frac {\ell (g ^ {n})}{n}
\]

存在，称为 g 的稳定换位子长度。

\begin{enumerate}
\def\labelenumi{\alph{enumi})}
\setcounter{enumi}{2}
\tightlist
\item
  设 S 为一个集合，\(\mathrm{G} = \mathrm{D}(\mathrm{F}(\mathrm{S}))\) 为由 S 生成的自由群的导出子群。群 G 等于其导出子群，且对每个 \(g \in G\)，\(g \neq e\)，有 \(\operatorname{lsc}(g) \geqslant 1/2\)。存在 \(g \in G\) 使 \(\operatorname{lsc}(g) = 1/2\)。
\end{enumerate}

\begin{enumerate}
\def\labelenumi{\arabic{enumi})}
\setcounter{enumi}{3}
\tightlist
\item
  设 \(n \geqslant 1\) 为整数。用 \(S(n)\) 表示满足下列条件的最大整数 \(k \geqslant 0\)：对每个集合 \(X\)（由 \(n\) 个实数组成），存在子集 \(Y \subset X\)（基数为 \(k\)），使方程 \(x + y = z\) 在 \(Y\) 中无解。于是 \(S(n) \leqslant n\)。
\end{enumerate}

\begin{enumerate}
\def\labelenumi{\alph{enumi})}
\tightlist
\item
  证明极限
\end{enumerate}

\[
\sigma = \lim _ {n \to + \infty} \frac {\mathrm{S} (n)}{n}
\]

存在。

\begin{enumerate}
\def\labelenumi{\alph{enumi})}
\setcounter{enumi}{1}
\item
  证明 \(\sigma \leqslant 2 / 5\)。
\item
  证明 \(\sigma \geqslant 1 / 3\)。
\end{enumerate}

（可以证明 \(\sigma = 1/3\)，参见 S. EBERHARD, B. GREEN, F. MANNERS, Sets of integers with no large sum-free subset, Annals of Mathematics 180 (2014), 621--652。）

\begin{enumerate}
\def\labelenumi{\arabic{enumi})}
\setcounter{enumi}{4}
\tightlist
\item
  设 A 为交换环，\(n \geqslant 1\) 为整数。对于系数在 A 中的矩阵 \(M = (m_{i,j})\)（其型为 \((n, n)\)），我们把 M 的积和式称为元素
\end{enumerate}

\[
\mathrm{per(M)} = \sum_ {\sigma \in \mathrm{S} _ {n}} m _ {1, \sigma (1)} \dots m _ {n, \sigma (n)}
\]

该元素属于 A。

设 \(k \geqslant 1\) 为整数。用 \(X_{n,k}\) 表示由实系数的、型为 \((n, n)\) 且每列恰含 k 个等于 1 的系数（其余系数全为零）的矩阵 M 所成的集合。

证明极限

\[
\lim _ {n \to + \infty} \left(\inf _ {\mathrm{M} \in \mathrm{X} _ {n, k}} \operatorname{per} (\mathrm{M})\right) ^ {1 / n}
\]

存在。

\begin{enumerate}
\def\labelenumi{\arabic{enumi})}
\setcounter{enumi}{5}
\tightlist
\item
  设 \(k \geqslant 1\) 为整数。对每个整数 \(N \geqslant 1\)，用 \(S_k(N)\) 表示不包含 \(A \subset \{1, \ldots, N\}\) 个成等差数列的元素的集合 \(k\) 的最大基数；也就是说，对任意 \(n_0 \geqslant 0\)、\(h \geqslant 1\)、\(\ell \geqslant 1\)，若 \(\{n_0 + h, n_0 + 2h, \ldots, n_0 + \ell h\} \subset A\)，则 \(\ell < k\)。
\end{enumerate}

\begin{enumerate}
\def\labelenumi{\alph{enumi})}
\tightlist
\item
  证明极限
\end{enumerate}

\[
s _ {k} = \lim _ {\mathrm{N} \to + \infty} \frac {\mathrm{S} _ {k} (\mathrm{N})}{\mathrm{N}}
\]

存在。

\begin{enumerate}
\def\labelenumi{\alph{enumi})}
\setcounter{enumi}{1}
\tightlist
\item
  证明 \(s_{1}=s_{2}=0\)。可以证明（“Szemerédi 定理”，参见 E. SZEMERÉDI, On sets of integers containing no k elements in arithmetic progression, Acta Arithmetica XXVII, 1975, 199--245）：对所有 \(s_{k}=0\) 有 \(k\geqslant1\)；至于 k=3 的情形，参见 II，第 270 页，习题 23。
\end{enumerate}

\begin{enumerate}
\def\labelenumi{\arabic{enumi})}
\setcounter{enumi}{6}
\tightlist
\item
  设 A 为赋范代数。
\end{enumerate}

\begin{enumerate}
\def\labelenumi{\alph{enumi})}
\tightlist
\item
  设 S 为 A 的非空有界子集。证明极限
\end{enumerate}

\[
\varrho (\mathrm{S}) = \lim _ {n \to + \infty} \sup _ {x \in \mathrm{S} ^ {n}} \| x \| ^ {1 / n}
\]

存在，并且当 S 仅由单个元素 x 组成时，有 \(\varrho(\mathrm{S}) = \varrho(x)\)。我们称 \(\varrho(\mathrm{S})\) 为子集 S 的谱半径。

\begin{enumerate}
\def\labelenumi{\alph{enumi})}
\setcounter{enumi}{1}
\tightlist
\item
  设 N 为 A 上与 A 的范数等价、且使赋以 p 之后的 A 仍为赋范代数的范数 p 全体所成的集合。设 X 为 A 的一个子集。为存在 \(p \in N\) 使得对所有 \(p(x) \leqslant 1\) 都有 \(x \in X\)，充要条件是存在实数 \(c \geqslant 0\) 使得
\end{enumerate}

\[
\sup _ {n \geqslant 1} \sup _ {x \in X ^ {n}} \| x \| \leqslant c.
\]

\begin{enumerate}
\def\labelenumi{\alph{enumi})}
\setcounter{enumi}{2}
\tightlist
\item
  证明
\end{enumerate}

\[
\varrho (\mathrm{S}) = \inf _ {p \in \mathrm{N}} \sup _ {x \in \mathrm{S}} p (x).
\]

\begin{enumerate}
\def\labelenumi{\arabic{enumi})}
\setcounter{enumi}{7}
\item
  设 A 为 \(\mathbf{R}\) 上在无穷远处趋于 0 的复值连续函数所成的代数，赋以范数 \(\|f\| = \sup_{t\in\mathbf{R}} |f(t)|\)。于是 \(\widetilde{\mathrm{A}}\) 就恒同于 \(\mathbf{R}\) 上在无穷远处趋于有限极限的复值连续函数所成的代数。在 \(\widetilde{\mathrm{A}}\) 上，范数 \(\|g\| = \sup_{t\in\mathbf{R}} |g(t)|\) 与 I，第 15 页，第 1 节中定义的范数不同。
\item
  设 A 为赋范代数，\(b = \inf_{x \neq 0} \|x^{2}\|/\|x\|^{2}\)。则
\end{enumerate}

\[
b \leqslant \inf _ {x \neq 0} \frac {\varrho (x)}{\| x \|} \leqslant \sqrt {b}.
\]

\begin{enumerate}
\def\labelenumi{\arabic{enumi})}
\setcounter{enumi}{9}
\item
  设 A 为赋范代数，\(x, y \in A\)。我们有 \(\varrho(xy) = \varrho(yx)\)。
\item
  设 H 为可数型 Hilbert 空间，\((f_{n})_{n\geqslant1}\) 为 H 的一个标准正交基。设 A 为赋范代数 \(\mathcal{L}(\mathrm{H})\)。设 \(x\in A\) 为由 \(x(f_{i})=2^{-i}f_{i+1}\) 所定义的元素。证明 x 是拟幂零元，但不是幂零元。
\end{enumerate}

\begin{enumerate}
\def\labelenumi{¶ \arabic{enumi})}
\setcounter{enumi}{11}
\tightlist
\item
  设 H 为可分 Hilbert 空间，\((f_n)_{n \geqslant 1}\) 为 H 的一个标准正交基。按下述方式定义诸数 \((\alpha_m)_{m \geqslant 1}\)：若 \(\alpha_m = e^{-k}\) 是 \(m\) 与一个奇数之积，则令 \(2^k\)。
\end{enumerate}

\begin{enumerate}
\def\labelenumi{\alph{enumi})}
\tightlist
\item
  设 \(x \in \mathcal{L}(\mathrm{H})\) 由 \(x(f_m) = \alpha_m f_{m+1}\)（对所有 \(m \geqslant 1\)）定义。证明 \(x\) 不是拟幂零元。（注意
\end{enumerate}

\[
\left\| x ^ {n} \right\| = \sup _ {m} \bigl (\alpha_ {m} \alpha_ {m + 1} \dots \alpha_ {m + n - 1} \bigr),
\]

并估计 \(\alpha_{1}\alpha_{2}\ldots\alpha_{2^{t}-1}\)。）

\begin{enumerate}
\def\labelenumi{\alph{enumi})}
\setcounter{enumi}{1}
\tightlist
\item
  设 \(x_{k} \in \mathcal{L}(\mathrm{H})\) 由下式定义：当 m 是 \(x_{k}(f_{m}) = 0\) 与一个奇数之积时 \(2^{k}\)，否则 \(x_{k}(f_{m}) = \alpha_{m} f_{m+1}\)。证明 \(x_{k}\) 是幂零元，且当 k 趋于 \((x_{k})\) 时 \(+\infty\) 收敛于 x。
\end{enumerate}

\begin{enumerate}
\def\labelenumi{\arabic{enumi})}
\setcounter{enumi}{12}
\item
  设 A 为幺 Banach 代数。设 \(x \in A\) 满足 \(\|x - 1\| < 1\)。则存在 \(y \in A\) 使 \(y^{2} = x\)。
\item
  设 A 为复幺赋范代数。若对 A 的每个可逆元 x 都有 \(\|x^{-1}\| = \|x\|^{-1}\)，则 A = C · 1。（可归结为 A 完备的情形。设 \(x \in A\) 与 C · 1 的距离为 \(\alpha > 0\)。若 \(\lambda_{0} \in C\) 使 \(x - \lambda_{0}\) 可逆，则当 \(x - \lambda\) 时 \(|\lambda - \lambda_{0}| < \alpha\) 也必可逆。由此推出 x 的谱将是空集。）
\item
  设 A 为复幺赋范代数，其中仅有的左拓扑零因子是 0，仅有的右拓扑零因子也是 0。则 \(A = C \cdot 1\)。（若 \(x \in A\)，且 \(\lambda\) 属于 \(\mathrm{Sp}_{\widehat{\mathrm{A}}}(\boldsymbol{x})\) 的边界，则 \(x - \lambda\) 是 \(\widehat{A}\) 中的拓扑零因子，因而是 A 中的拓扑零因子，于是 \(x = \lambda\)。）
\item
  设 A 为复幺 Banach 代数。设 \(a\) 为 A 的一个元素。我们把 a 的奇异谱称为满足下列条件的 \(\lambda \in \mathbf{C}\) 所成的集合：\(a - \lambda e\) 在 A 中既是左拓扑零因子又是右拓扑零因子。
\end{enumerate}

\begin{enumerate}
\def\labelenumi{\alph{enumi})}
\item
  证明 a 的奇异谱包含于 a 的谱中，并且包含 a 的谱的边界。
\item
  设 \(f \in C[X]\) 为多项式。a 的奇异谱在 f 下的像就是 \(f(a)\) 的奇异谱。
\end{enumerate}

\begin{enumerate}
\def\labelenumi{\arabic{enumi})}
\setcounter{enumi}{16}
\tightlist
\item
  对每个 \(x \in A\)，我们令：
\end{enumerate}

\[
\lambda (x) = \inf _ {y \neq 0} \frac {\| x y \|}{\| y \|} \quad \lambda^ {\prime} (x) = \inf _ {y \neq 0} \frac {\| y x \|}{\| y \|}.
\]

\begin{enumerate}
\def\labelenumi{\alph{enumi})}
\tightlist
\item
我们有
\end{enumerate}

\[
| \lambda (x) - \lambda (y) | \leqslant \| x - y \|, \quad | \lambda^ {\prime} (x) - \lambda^ {\prime} (y) | \leqslant \| x - y \|,
\]

\[
\lambda (x) \lambda (y) \leqslant \lambda (x y) \leqslant \| x \| \lambda (y), \quad \lambda^ {\prime} (x) \lambda^ {\prime} (y) \leqslant \lambda^ {\prime} (x y) \leqslant \lambda^ {\prime} (x) \| y \|.
\]

\begin{enumerate}
\def\labelenumi{\alph{enumi})}
\setcounter{enumi}{1}
\tightlist
\item
  A 中的左（相应地，右）拓扑零因子所成的集合是闭集。
\end{enumerate}

\begin{enumerate}
\def\labelenumi{\arabic{enumi})}
\setcounter{enumi}{17}
\item
  设 A 为幺 Banach 代数。由既不可逆又不是拓扑零因子的 \(x\) 所成的集合在 A 中是开集。
\item
  设 X 为复 Banach 空间，\(x \in A = \mathcal{L}(X)\)。
\end{enumerate}

\begin{enumerate}
\def\labelenumi{\alph{enumi})}
\tightlist
\item
  若 x 不可逆，则 x 是左拓扑零因子或右拓扑零因子。（依次考虑下列各种情形：）
\end{enumerate}

\(1^{\circ}\) ) \(x\) 不是单射；

\(2^{\circ})\overline{x(\mathrm{X})}\neq \mathrm{X};\)

\(3^{\circ}) x \text{ 是单射的， } x(\mathrm{X}) \text{ 在 X 中稠密且不同于 X。}\)

\begin{enumerate}
\def\labelenumi{\alph{enumi})}
\setcounter{enumi}{1}
\tightlist
\item
  若 x 把 X 双连续地映到 X 的一个异于 X 的闭向量子空间上，或者 x 非单射而其像为 X，则 x 位于 A 的不可逆元素所成集合的内部。
\end{enumerate}

\begin{enumerate}
\def\labelenumi{\arabic{enumi})}
\setcounter{enumi}{19}
\item
  在 I，第 20 页，例 9 的代数 A 中，恒等函数 z 既不可逆，也不是拓扑零因子。
\item
  设 A 为幺赋范代数。设 B 为赋范代数 \(\mathcal{L}(\mathrm{A})\)。于是 A 在态射 \(x \mapsto \gamma_x\) 下的像是 B 的一个全子代数。因此，若 \(x \in \mathrm{A}\)，则 \(\mathrm{Sp}_{\mathrm{A}}(x) = \mathrm{Sp}_{\mathrm{B}}(\gamma_x)\)。
\end{enumerate}

\begin{enumerate}
\def\labelenumi{¶ \arabic{enumi})}
\setcounter{enumi}{21}
\tightlist
\item
  设 A 为 Banach 代数。
\end{enumerate}

\begin{enumerate}
\def\labelenumi{\alph{enumi})}
\tightlist
\item
  设 S 为由 A 的元素组成的有界序列全体，赋以范数 \(\|(\boldsymbol{x}_{n})\| = \sup \|x_{n}\|\)。设 R 为 \((x_{n}) \in \mathrm{S}\) 中使得 \(\|x_{n}\|\) 趋于 0 的元素所成的集合。于是 S 是 Banach 代数，而 R 是 S 的一个闭双边理想。
\end{enumerate}

设 \(S' = \text{S/R}\)，并设 \(\varphi: \text{S} \to \text{S}'\) 为典范同态。

\begin{enumerate}
\def\labelenumi{\alph{enumi})}
\setcounter{enumi}{1}
\item
  映射 \(\theta\)（由 \(\theta(x) = \varphi(x, x, x, \ldots)\) 定义，从 A 到 S' 中）是 A 到 S' 的一个子代数上的等距同构。S' 中的每个左拓扑零因子都是左零因子。为使 \(x \in A\) 是左拓扑零因子，充要条件是 \(\theta(x)\) 是 S' 中的左零因子。
\item
  假设 A 有单位元。证明：存在一个 Banach 代数 B 以及 A 到 B 的一个子代数上的等距同构，使得 A 的元素不可逆当且仅当它在 B 中的像是左零因子或右零因子。（利用 a) 以及习题 19 与 21。）
\end{enumerate}

\begin{enumerate}
\def\labelenumi{\arabic{enumi})}
\setcounter{enumi}{22}
\tightlist
\item
  设 A 为赋范代数，R 为其根。
\end{enumerate}

\begin{enumerate}
\def\labelenumi{\alph{enumi})}
\item
  若 \(x \in R\)，则 x 是拟幂零元。（我们有 \(\mathrm{Sp}_{\mathrm{A}}^{\prime}(x) = \{0\}\)，从而更有 \(\mathrm{Sp}_{\widehat{\mathrm{A}}}^{\prime}(x) = \{0\}\)。）
\item
  假设 A 是 Banach 代数。设 I 为 A 的一个左理想，其所有元素都是拟幂零元。则 \(I \subset R\)（可以假设 A 有单位元。若 \(x \in R\) 且 \(a \in A\)，则有 \(\mathrm{Sp}_{\mathrm{A}}(ax) = \{0\}\)，因而 1 - ax 可逆；从而 \(x \in R\)。
\end{enumerate}

\begin{enumerate}
\def\labelenumi{\arabic{enumi})}
\setcounter{enumi}{23}
\item
  设 A 为复 Banach 代数，B 为无根的复代数，\(\varphi\) 为从 A 到 B 上的态射。\(\varphi\) 的核是闭的。
\item
  设 A 为复 Banach 代数，\(\pi\) 为 A 在复向量空间 X 中的一个非零不可约表示。设 \(\xi_{0}\) 为 X 的非零元素。
\end{enumerate}

\begin{enumerate}
\def\labelenumi{\alph{enumi})}
\item
  \(\xi_0\) 的零化子 I 是 A 的一个正则极大左理想；它是闭的。
\item
  映射 \(x \mapsto x\xi_{0}\) 通过过渡到商，定义了 A-模 A/I 到 A-模 X 上的一个同构 \(\varphi\)。
\item
  若经 \(\varphi\) 把 Banach 空间 A/I 的范数转移过来，则 X 成为 Banach 空间，且对所有 \(\| \pi (x)\| \leqslant \| x\|\) 有 \(x\in \mathrm{A}\)。
\item
  若 A 是本原的（即 \(\{0\}\) 是本原理想，参见 I，第 11 页的定义，另参见 A，VIII，第 433 页，习题 6），则依照 A 有或没有单位元而分别有 \(A = \{0\}\) 或 \(C \cdot 1\)。（利用前面的结果以及 I，第 26 页的推论 4。）
\end{enumerate}

\begin{enumerate}
\def\labelenumi{¶ \arabic{enumi})}
\setcounter{enumi}{25}
\tightlist
\item
  设 A 为 Banach 代数，R 为其根。
\end{enumerate}

\begin{enumerate}
\def\labelenumi{\alph{enumi})}
\tightlist
\item
  证明：若 \(r \in R\)，则级数
\end{enumerate}

\[
- \frac {1}{2} \sum_ {k = 1} ^ {\infty} \binom{1 / 2}{k} (- 4 r) ^ {k}
\]

收敛于满足 \(x \in R\) 的元素 \(x^{2} - x + r = 0\)。

\begin{enumerate}
\def\labelenumi{\alph{enumi})}
\setcounter{enumi}{1}
\tightlist
\item
  设 u 为 A 的一个元素，它模 R 的类是幂等的。则存在 A 中与 u 模 R 同余的一个幂等元。（可以假设 A 有单位元。设 \(q = u - u^{2} \in R\)，并设 \(x \in R\) 为借助 a) 构造的方程 \(x^{2} - x - q(1 - 4q)^{-1} = 0\) 的一个解。那么 \(u - (2u - 1)x\) 即为所求。）
\end{enumerate}

\begin{enumerate}
\def\labelenumi{\arabic{enumi})}
\setcounter{enumi}{26}
\item
  设 A 为复幺 Banach 代数，\(x \in A\)，并设 \(\zeta\) 为 \(\mathrm{Sp}_{\mathrm{A}}(x)\) 的一个边界点。x 的预解式不能延拓为在 \(\zeta\) 处连续的函数。
\item
  设 A 为复幺 Banach 代数，\(\Delta\) 为 C 的开子集，\(\lambda \mapsto \mathrm{R}(\lambda)\) 为从 \(\Delta\) 到 A 中的映射，使得
\end{enumerate}

\[
\mathrm{R} (\lambda) - \mathrm{R} (\mu) = - (\lambda - \mu) \mathrm{R} (\lambda) \mathrm{R} (\mu)\tag{5}
\]

对所有 \(\lambda,\mu\in\Delta\) 成立。

\begin{enumerate}
\def\labelenumi{\alph{enumi})}
\tightlist
\item
  设 \(\lambda_0\in \Delta\)。对每个满足 \(\lambda \in \Delta\) 的 \(|\lambda -\lambda_0|\cdot \| \mathrm{R}(\lambda)\| <  1\)，我们有：
\end{enumerate}

\[
\mathrm{R} (\lambda) = \sum_ {n = 0} ^ {\infty} (\lambda_ {0} - \lambda) ^ {n} \mathrm{R} (\lambda_ {0}) ^ {n + 1}.
\]

\begin{enumerate}
\def\labelenumi{\alph{enumi})}
\setcounter{enumi}{1}
\tightlist
\item
  对每个 \(\lambda \in \Delta\) 和每个整数 \(n\geqslant 0\)，有
\end{enumerate}

\[
\frac {\partial^ {n}}{\partial \lambda^ {n}} \mathrm{R} (\lambda) = (- 1) ^ {n} n! \mathrm{R} (\lambda) ^ {n + 1}.
\]

\begin{enumerate}
\def\labelenumi{\alph{enumi})}
\setcounter{enumi}{2}
\tightlist
\item
  若 R 在无穷远处全纯，则存在 \(z, j, x \in \mathrm{A}\) 使得 \(z^2 = 0\)、\(j^2 = j\)、\(zj = jz = 0\)、\(x \in j\mathrm{Aj}\)，且当 \(\mathrm{R}(\lambda) = z + j\mathrm{R}(\lambda, x)\) 充分大时 \(|\lambda|\)。（当 \(\mathrm{R}(\lambda) = \sum_{n=0}^{\infty} c_n \lambda^{-n}\) 时写出 \(|\lambda| > \lambda_0\)，并表出 R 满足 (5)。可令 \(c_0 = z\)、\(c_1 = j\)、\(c_2 = x\)。）
\end{enumerate}

\begin{enumerate}
\def\labelenumi{\alph{enumi})}
\setcounter{enumi}{3}
\tightlist
\item
  为存在 \(a \in A\) 使 \(R\) 是 \(\Delta\) 的预解式 \(\lambda \mapsto R(a, \lambda)\) 在 \(a\) 上的限制，充要条件是存在 \(\lambda_0 \in \Delta\) 使 \(R(\lambda_0)\) 可逆。
\end{enumerate}

\begin{enumerate}
\def\labelenumi{\alph{enumi})}
\setcounter{enumi}{4}
\tightlist
\item
  对每个 \(x \in A\)，函数 \(\lambda \mapsto \sum_{n=0}^{\infty} (\lambda_0 - \lambda)^n x^{n+1}\) 在条件 \(|\lambda - \lambda_0| \cdot \|x\| < 1\) 下满足 (5)。
\end{enumerate}

\begin{enumerate}
\def\labelenumi{\arabic{enumi})}
\setcounter{enumi}{28}
\tightlist
\item
  \begin{enumerate}
  \def\labelenumii{\alph{enumii})}
  \tightlist
  \item
    设 E 为复 Banach 空间，\(f: \mathbf{C} \to \mathbf{E}\) 为整函数，\(\theta \mapsto \mathrm{P}(\theta) = \sum_{\mu=k}^{m} c_{\mu} e^{\mu i\theta}\) 为三角多项式（\(c_k \neq 0\)）。假设当 \(|\mathrm{P}(\theta)| \cdot \|f(re^{i\theta})\| \leqslant \mathrm{Mr}^{\alpha}\) 时 \(r \geqslant r_0\)（M、\(r_0\)、\(\alpha\) 为常数，均 \(\geqslant 0\)）。则 \(f\) 是次数 \(\leqslant \alpha\) 的多项式。（设 \(f(\zeta) = \sum_{\nu=0}^{\infty} a_{\nu} \zeta^{\nu}\)，并令
  \end{enumerate}
\end{enumerate}

\[
\mathrm{J} (r, p) = r ^ {- p} \int_ {0} ^ {2 \pi} \mathrm{P} (\theta) f (r e ^ {i \theta}) e ^ {- (k + p) i \theta} d \theta .
\]

对 \(\mathrm{J}(r,p)\)，用两种不同的方式计算当 \(r\to+\infty\) 时 \(p>\alpha.\) 的极限。

\begin{enumerate}
\def\labelenumi{\alph{enumi})}
\setcounter{enumi}{1}
\tightlist
\item
  沿用 \(a\) 的记号，若 \(\alpha\) 是整数，且对每个 \(\theta \in \mathbf{R}\) 有
\end{enumerate}

\[
\lim _ {r \to + \infty} r ^ {- \alpha} | \mathrm{P} (\theta) | \| f (r e ^ {i \theta}) = 0,
\]

则 \(f\) 是次数 \(< \alpha\) 的多项式。

\begin{enumerate}
\def\labelenumi{\alph{enumi})}
\setcounter{enumi}{2}
\tightlist
\item
  设 A 为复幺 Banach 代数，x 为 A 的一个拟幂零元。假设对某个整数 n \textgreater{} 0，且对每个 \(\theta \in R\)，有
\end{enumerate}

\[
\lim _ {r \to + \infty} r ^ {- n} | \cos^ {n} \theta | \cdot \| (1 - r e ^ {i \theta} x) ^ {- 1} \| = 0.
\]

则 \(x^n = 0\)。（利用 \(b\)。）

\begin{enumerate}
\def\labelenumi{¶ \arabic{enumi})}
\setcounter{enumi}{29}
\tightlist
\item
  a) 设 E 为复 Banach 空间，\(x: \mathbf{C}=\{1\}\to\mathrm{E}\) 为 \(1/(\zeta-1)\) 的整函数。用 \(x(\zeta)=\sum_{n=0}^{\infty}a_{n}\zeta^{n}\)、\(x(\zeta)=\sum_{n=0}^{\infty}b_{n}\zeta^{-n}\) 分别表示 \(x\) 在 \(|\zeta|<1\)、\(|\zeta|>1\) 中的展开式。假设存在 \(\alpha>0\) 使得 \(\|a_{n}\|=o(n^{\alpha})\)、\(\|b_{n}\|=o(n^{\alpha})\)。则 \(x(\zeta)\) 是 \(1 / (\zeta - 1)\) 的次数 \(< \alpha + 1\) 的多项式。（设 \(\varepsilon > 0\)。我们有：当 \(\|x(\zeta)\| \leqslant \varepsilon(1 - r)^{-1 - \alpha}\) 时 \(r_{\varepsilon} \leqslant r = |\zeta| < 1\)；当 \(\|x(\zeta)\| \leqslant \varepsilon(1 - r^{-1})^{-1 - \alpha}\) 时 \(1 < r \leqslant r_{\varepsilon}'\)。令 \(1 - \zeta = (\omega + \frac{1}{2})^{-1}\)、\(\omega = \mathrm{Re}^{i\theta}\)，证明：对 \(0 \leqslant r < 1\) 与 \(\mathrm{R} \geqslant 1\)，有
\end{enumerate}

\[
(1 - r) ^ {- 1} \leqslant \frac {9}{4} \frac {\mathrm{R}}{\cos \theta}
\]

并且，对 1 \textless{} r 与 R ≥ 1，

\[
(1 - r ^ {- 1}) ^ {- 1} \leqslant \frac {9}{4} \frac {\mathrm{R}}{\cos \theta}.
\]

令 \(y(\omega) = x(\zeta)\)，并证明当 \(R^{-\alpha - 1}|\cos \theta|^{\alpha + 1}\| y(Re^{i\theta})\|\) 趋于 \(R\) 时 \(+\infty\) 关于 \(\theta\) 一致地趋于 0。然后应用习题 29 的 b)。

设 A 为复幺 Banach 代数，\(q \in A\) 为拟幂零元，且 \(x = 1 + q\)。

\begin{enumerate}
\def\labelenumi{\alph{enumi})}
\setcounter{enumi}{1}
\item
  为使 \(q^{N}=0\)，充要条件是当 n 趋于 \(\|x^{\pm n}\|=o(n^{N})\) 时 \(+\infty\)。（为看出条件是充分的，应用 a) 证明 \(\mathrm{R}(\lambda,x)\) 是 \(1/(\lambda-1)\) 的次数 \(\leqslant N\) 的多项式。另一方面，\(\mathrm{R}(\lambda,x)=\sum_{n=0}^{\infty}q^{n}(\lambda-1)^{-n-1}\)。）
\item
  设 R 为 A 的根，G 为 A 的可逆元群的一个子群。若对每个 \(x \in G\)，当 \(\|x^{n}\| = o(n)\) 时有 \(\|x^{-n}\| = o(n)\)、\(n \to +\infty\)，则典范映射 \(A \to A/R\) 在 G 上的限制是单射。
\end{enumerate}

\begin{enumerate}
\def\labelenumi{\arabic{enumi})}
\setcounter{enumi}{30}
\tightlist
\item
  \begin{enumerate}
  \def\labelenumii{\alph{enumii})}
  \tightlist
  \item
    设 A 为赋范代数，x、y 为 A 中满足 xy - yx = y 的两个元素。证明：当 \(y^{n} = 0\) 时 \(n > 2\|x\|\)。（利用公式 \(xy^{n} - y^{n}x = ny^{n}\)。）
  \end{enumerate}
\end{enumerate}

\begin{enumerate}
\def\labelenumi{\alph{enumi})}
\setcounter{enumi}{1}
\item
  设 L 为 C 上的有限维 Lie 代数。证明：若 L 不是幂零的，则它含有两个非零元素 x、y，使得 \([x, y] = y\)。（利用存在 \(x \in L\) 使 \(\operatorname{ad}(x)\) 不是幂零的这一事实。）
\item
  设 L 为实或复的有限维 Lie 代数，且 L 的包络代数具有一个与其代数结构相容的范数。证明 L 是幂零的。（利用 a) 与 b)。）
\end{enumerate}

\begin{enumerate}
\def\labelenumi{\alph{enumi})}
\setcounter{enumi}{3}
\tightlist
\item
  设 V 为实或复的有限维向量空间，\(\mathfrak{L}\) 为 \(\mathfrak{gl}(\mathrm{V})\) 的一个由幂零自同态组成的 Lie 子代数，\(x_{1},\ldots ,x_{n}\) 为 \(\mathfrak{L}\) 中的元素。设 \(\alpha >0\)。证明：V 上存在一个 Hilbert 空间结构，使得对 \(\| x_i\| \leqslant \alpha\) 有 \(1\leqslant i\leqslant n\)。
\end{enumerate}

\begin{enumerate}
\def\labelenumi{\alph{enumi})}
\setcounter{enumi}{4}
\tightlist
\item
  设 \(\mathfrak{L}\) 为实或复的有限维幂零 Lie 代数，U 为其包络代数。证明：U 上存在一个与其代数结构相容的范数。（利用 LIE，I，§3，习题 5 与 §7，习题 3 的方法，证明存在一族表示 \((\pi_{\lambda})\)，它们是 \(\mathfrak{L}\) 在有限维空间 \(V_{\lambda}\) 中的表示，具有下列性质：
\end{enumerate}

对所有 \(u \in U\)，存在 \(\lambda\) 使 \(\pi_{\lambda}(u) \neq 0\)；并且对每个 \(\pi_{\lambda}(\mathfrak{L})\)，\(\lambda\) 都由幂零自同态组成。给每个 \(V_{\lambda}\) 赋以由 d) 提供的 Hilbert 结构，便在 Hilbert 空间 V（诸 \(V_{\lambda}\) 的 Hilbert 和）中得到 L 的一个由连续自同态给出的表示 \(\pi\)，从而得到一个单射态射 \(\varphi\)，它把 U 映入 \(\mathcal{L}(V)\)。令 \(\|u\|_{U} = \|\varphi(u)\|\)（对所有 \(u \in U\)）。

\begin{enumerate}
\def\labelenumi{\alph{enumi})}
\setcounter{enumi}{5}
\tightlist
\item
  设 A 为复幺赋范代数，x、y 为 A 的元素。若 \(A \neq \{0\}\)，则 xy - yx ≠ 1。（第一个证明：我们有 \(\mathrm{Sp}_{\mathrm{A}}^{\prime}(xy) = \mathrm{Sp}_{\mathrm{A}}^{\prime}(yx)\)；若 xy = yx + 1，则 \(\mathrm{Sp}_{\mathrm{A}}(xy)\) 可由 \(\mathrm{Sp}_{\mathrm{A}}(yx)\) 经平移 \(z \mapsto z + 1\) 得到；由此推出 \(\mathrm{Sp}_{\mathrm{A}}(xy)\) 无界，而这是荒谬的。第二个证明：若 \([x, y] = 1\)，则有 \([xy, y] = y\)，从而由 a) 知对充分大的 n 有 \(y^{n} = 0\)；再由 \([x, y^{p}] = py^{p-1}\) 推出 y = 0。）
\end{enumerate}

\begin{enumerate}
\def\labelenumi{\arabic{enumi})}
\setcounter{enumi}{31}
\item
  设 A 为赋范代数，\(x \in A\)。若当 n 趋于 \(x^{n}\) 时 \(+\infty\) 趋于 0，则称 x 是拓扑幂零的。证明：x 是拓扑幂零的当且仅当 \(\varrho(x) < 1\)。
\item
  设 A 为 C 上的交换幺 Banach 代数。假设 A 的每个闭理想 I 都是有限生成的（即存在 \(n \in N\) 与 \(x_{1}, \ldots, x_{n} \in A\)，使得 \(I = Ax_{1} + \cdots + Ax_{n}\)）。
\end{enumerate}

\begin{enumerate}
\def\labelenumi{\alph{enumi})}
\tightlist
\item
  A 的每个理想 I 都是闭的。（写 \(\bar{\mathrm{I}} = \mathrm{A}x_{1} + \dots +\mathrm{A}x_{n}\)，其中 \(x_{1},\ldots ,x_{n}\) 属于 A。设 \((x_{ip})_{p\geqslant 1}\) 是 I 中趋于 \(x_{i}\) 的元素序列。对每个 \((y_{1},\ldots ,y_{n})\in \mathrm{A}^{n}\) 与每个 \(p\geqslant 1\)，令
\end{enumerate}

\[
\begin{array}{c} \varphi (y _ {1}, \ldots , y _ {n}) = x _ {1} y _ {1} + \dots + x _ {n} y _ {n} \in \bar {\mathrm{I}} \\ \varphi_ {p} (y _ {1}, \ldots , y _ {n}) = x _ {i p} y _ {1} + \dots + x _ {n p} y _ {n} \in \mathrm{I}. \end{array}
\]

我们有 \(\varphi\in\mathcal{L}(A^{n},\bar{I})\)、\(\varphi_{p}\in\mathcal{L}(A^{n},\bar{I})\)，且 \(\|\varphi-\varphi_{p}\|\) 随 \(p\) 趋于 \(+\infty\) 而趋于 0，并且 \(\varphi\) 是满射的。借助 \(\varphi_{p}\) 当 \(p\) 充分大时的满射性（考虑 \(\varphi\) 与 \(\varphi_{p}\) 的转置），推出 I = \(\bar{I}\)。

\begin{enumerate}
\def\labelenumi{\alph{enumi})}
\setcounter{enumi}{1}
\item
  若 A 是整环，则 \(A = C \cdot 1\)。（利用 a) 与习题 15。）
\item
  若 A 中唯一的幂零元是 0，则存在整数 \(n \geqslant 0\) 使 A 同构于 \(\mathbf{C}^n\)。（由于 A 是 Noether 的，\(\{0\}\) 是素理想 \(\mathfrak{p}_1, \ldots, \mathfrak{p}_m\) 的交；由 \(a\)) 知它们都是闭的；而代数 \(\mathrm{A}/\mathfrak{p}_i\) 同构于 \(\mathbf{C}\)（据 \(b\)）。）
\end{enumerate}

\begin{enumerate}
\def\labelenumi{\alph{enumi})}
\setcounter{enumi}{3}
\tightlist
\item
  我们有 \(\dim_{\mathbf{C}}\mathrm{A} < +\infty\)。（设 N 为 A 的幂零元全体，它是 A 的一个（闭）理想。于是由 \(\dim (\mathrm{A / N}) < +\infty\)) 有 \(c\)。由于 N 是有限生成理想，故对充分大的 \(\mathrm{N}^n = 0\) 有 \(n\)。最后，每个 \(\mathrm{N}^i /\mathrm{N}^{i + 1}\) 都是 A/N 上的有限生成模，故 \(\dim (\mathrm{N}^i /\mathrm{N}^{i + 1}) < +\infty\)。）
\end{enumerate}

\begin{enumerate}
\def\labelenumi{\arabic{enumi})}
\setcounter{enumi}{33}
\tightlist
\item
  设 A 为 C 上的幺代数，带有一个分离的局部凸拓扑，使得 A 中的乘法分别连续。
\end{enumerate}

。称元素 \(a \in A\) 为正则的，如果存在 \(r \geqslant 0\)，使得 \(a - \lambda\) 对 \(|\lambda| \geqslant r\) 可逆，且使得 \((a - \lambda)^{-1}\)（对 \(|\lambda| \geqslant r\)）所成的集合有界。

\begin{enumerate}
\def\labelenumi{\alph{enumi})}
\item
  若 a 正则，则 \(\lambda(a-\lambda)^{-1}\) 对 \(|\lambda|\geqslant r\) 有界。
\item
  若映射 \(x \mapsto x^{-1}\) 在 1 的一个邻域中有定义且在点 1 处连续，则 A 的每个元素都是正则的。
\item
  对每个 \(a \in A\)，设 \(U_{a}\) 为使 \(\lambda \in C\)，且 \((a - \mu)^{-1}\)（对 \(\mu\) 取遍 \(\lambda\) 的一个邻域）存在且有界的元素所成的集合。设 \(S_{a} = C - U_{a}\)。则 \(S_{a}\) 是闭的；且若 a 正则，则 \(S_{a}\) 是紧的且非空（假定 \(A \neq \{0\}\)）。
\item
  设 \(a \in A\)。函数 \(\lambda \mapsto (a - \lambda)^{-1}\) 在 \(U_{a}\) 中全纯。
\item
  设 \(a \in A\)。为使 \(\lambda \in U_{a}\)，充要条件是 \(a - \lambda\) 有正则逆。
\item
  若 A 的每个元素都正则，且 A 是域，则 \(A = C \cdot 1\)。
\end{enumerate}

\begin{enumerate}
\def\labelenumi{\arabic{enumi})}
\setcounter{enumi}{34}
\tightlist
\item
  设 A 为 R 上的一个代数且是域，带有一个分离的局部凸拓扑，使得 A 中的乘法连续，且映射 \(x \mapsto x^{-1}\) 在 1 处连续。则 A 与 R、与 C 或与四元数域 H 是 R-同构的。（若 A 交换且存在 \(u \in A\) 满足 \(u^{2} = -1\)，则给 A 一个 C 上的代数结构，并应用习题 34 的 f)。）若 A 交换且不存在元素 \(u \in A\) 满足 \(u^{2} = -1\)，则对域 \(A \otimes_{R} C\) 应用习题 34 的 f)。若 A 非交换，则按 AC，VI，§6，第 4 节，定理 1 的第三种情形论证。
\end{enumerate}

\begin{enumerate}
\def\labelenumi{¶ \arabic{enumi})}
\setcounter{enumi}{35}
\tightlist
\item
  设 A 为由复系数单变量有理函数到 {[}0,1{]} 上的限制所组成的 C 上代数。这个代数是一个域。
\end{enumerate}

\begin{enumerate}
\def\labelenumi{\alph{enumi})}
\item
  给 A 赋以依测度收敛的拓扑（INT，IV，§5，第 11 节）。则映射 \((x,y)\mapsto x+y\) 与 \((x,y)\mapsto xy\)（从 A × A 到 A 中）连续，且映射 \(x\mapsto x^{-1}\)（从 A* 到 A 中）连续。A 的拓扑不是局部凸的。
\item
  对 \(n \in N^{*}\)，定义序列 \((w_{r,n})_{r \in \mathbf{Z}}\) 如下：若 r \(w_{r,n} = (-r + 1)^{n(-r+1)}\)\textless{} 0，\(w_{0,n} = 1\)，且若 r \(w_{r,n} = (r + 1)^{-(r+1)/n}\)\textgreater{} 0。若 \(f \in A\)，令 \(p_n(f) = \sum_r w_{r,n} |a_r| < +\infty\)，其中 \(\sum a_r t^r\) 是 \(f(t)\) 在 0 处的 Laurent 展开式。于是 \(p_n\) 是一些半范数，它们在 A 上定义了一个可度量的局部凸拓扑。映射 \((x, y) \mapsto xy\)（从 A × A 到 A 中）连续。映射 \(x \mapsto x^{-1}\)（从 A* 到 A 中）不连续。代数 A 对此拓扑不完备。
\end{enumerate}

\begin{enumerate}
\def\labelenumi{\arabic{enumi})}
\setcounter{enumi}{36}
\tightlist
\item
  \begin{enumerate}
  \def\labelenumii{\alph{enumii})}
  \tightlist
  \item
    设 A 为 C 上带有局部凸拓扑的一个代数。下列条件等价：
  \end{enumerate}
\end{enumerate}

\begin{enumerate}
\def\labelenumi{(\roman{enumi})}
\item
  存在 0 的凸平衡邻域的基本组 \((\mathrm{U}_{i})\)，使得对一切 i 有 \(U_{i}U_{i} \subset U_{i}\)；
\item
  A 同构于一族赋范代数的乘积的一个子代数。
\item
  A 的拓扑可由一族半范数 \(p_{i}\) 定义，\(p_{i}(xy) \leqslant p_{i}(x)p_{i}(y)\) 对一切 \(x, y \in A\) 成立。
\end{enumerate}

若 A 满足这些条件，则称 A 是局部 m-凸的。

\begin{enumerate}
\def\labelenumi{\alph{enumi})}
\setcounter{enumi}{1}
\item
  设 A 为局部 m-凸代数。映射 \((x,y)\mapsto xy\)（从 A × A 到 A 中）连续。若 A 有单位元，且用 G 表示 A 的可逆元集合，则映射 \(x\mapsto x^{-1}\)（从 G 到 A 中）连续。
\item
  设 A 为幺局部 m-凸代数。A 的每个元素的谱非空。若 A 是域，则 \(A = C \cdot 1\)。
\item
  设 A 为完备的局部 m-凸代数。存在一个有向集 I、一族 Banach 代数 \((\mathrm{A}_{i})_{i\in\mathrm{I}}\)、以及对 \(\pi_{ij}\colon A_{j}\mapsto A_{i}\) 的连续态射 \(i\leqslant j\)，满足 \(\pi_{ij}\circ\pi_{jk}=\pi_{ik}\)，使得 A 同构于 \(\prod_{i}A_{i}\) 的由那些 \((x_{i})\)（它们满足 \(\pi_{ij}(x_{j})=x_{i}\)，对 \(i\leqslant j\)）组成的子代数。
\item
  代数 \(\mathcal{C}_{\mathbf{C}}(\mathbf{R})\)，赋以紧收敛拓扑，是局部 m-凸的。\(\mathcal{C}_{\mathbf{C}}(\mathbf{R})\) 的可逆元集合在 \(\mathcal{C}_{\mathbf{C}}(\mathbf{R})\) 中稠密；它不是开集。
\item
  设 D 为 C 的开子集，A 为 D 上复值全纯函数所成的代数，赋以紧收敛拓扑。则 A 是局部 m-凸的。A 的可逆元集合是闭的。
\item
  {[}0,1{]} 上无限次可微的复值函数所成的代数，赋以各阶导数一致收敛的拓扑，是局部 m-凸的。
\item
  设 A 为 {[}0,1{]} 上满足：对每个 p \(f \in \mathrm{L}^{p}([0,1])\)\textgreater{ 1 有 } 的复函数（的类）所成的代数。赋以由半范数族 \(f \mapsto \|f\|_{p}\)（p \textgreater{} 1）定义的拓扑，A 不是局部 m-凸的。映射 \((x,y) \mapsto xy\)（从 A × A 到 A 中）连续。
\end{enumerate}

\exercisesection{交换 Banach 代数}

在以下习题中，所考虑的代数除特别声明外均以 C 为基域。

\begin{enumerate}
\def\labelenumi{\arabic{enumi})}
\item
  设 A 为交换 Banach 代数，I 为 A 的一个极大理想。证明只有两种可能的情形：或者 I 是 A 的某个特征标的核，或者 I 是包含 \(A^{2}\) 的一个超平面。（利用 I，第 154 页的习题 6。）特别地，若 I 不是闭的，则 I 是 A 的包含 \(A^{2}\) 的稠密超平面。（为得到后一情形的一个例子，参见习题 3。）
\item
  设 A 为由趋于 0 的复数序列 \(x = (x_{n})_{n \geqslant 1}\) 组成的 Banach 代数，赋以范数 \(\|x\| = \sup_{i} |x_{i}|\)。设 I 为 A 中具有限支撑的元素所成的集合。则 I 是 A 的一个稠密理想，且不含于任何极大理想之中。（注意 \(A^{2} = A\)，并利用习题 1。）
\item
  设 \(\Lambda\) 为一个集合，\(p \in [1, +\infty[\)，并设 \(\mathbf{A} = \ell^p(\Lambda)\) 为由族 \(x = (x_{\lambda})_{\lambda \in \Lambda}\)（其中 \(x_{\lambda} \in \mathbf{C}\)）组成的集合，且
\end{enumerate}

\[
\| x \| _ {p} = \left(\sum_ {\lambda} | x _ {\lambda} | ^ {p}\right) ^ {1 / p} <   + \infty .
\]

\begin{enumerate}
\def\labelenumi{\alph{enumi})}
\item
  对按分量相加与相乘的运算以及范数 \(x \mapsto \| x \|_p\)，空间 A 是一个 Banach 代数。
\item
  我们有 \(A^{2} \neq A\)，且 \(A^{2}\) 在 A 中稠密。
\item
  A 的非零特征标所成的空间 \(\mathsf{X}(\mathsf{A})\) 典范地等同于赋有离散拓扑的 \(\Lambda\)。（利用 \(\ell^{p}(\Lambda)\) 的对偶是 \(\ell^{q}(\Lambda)\) 这一事实，其中 q 是 p 的共轭指数。）于是 Gelfand 变换成为恒等映射，它在 \(\Lambda\) 上在无穷远处趋于 0 的连续函数空间中的像是稠密的且不是闭的。
\end{enumerate}

\begin{enumerate}
\def\labelenumi{\arabic{enumi})}
\setcounter{enumi}{3}
\tightlist
\item
  设 \((\alpha_{n})_{n\geqslant 0}\) 是 \(\mathbf{R}_{+}^{*}\) 中的序列，满足 \(\alpha_0 = 1\)，对一切整数 \(\alpha_{m + n}\leqslant \alpha_m\alpha_n\) 与 \(m\) 有 \(n\)，且 \(\alpha_{n}^{1 / n}\to 0\)。设 A 为形式级数 \(x = \sum_{n = 0}^{\infty}x_{n}z^{n}\in \mathbf{C}[[z]]\) 中满足下列条件的那些所成的集合：
\end{enumerate}

\[
\| x \| = \sum_ {n = 0} ^ {\infty} | x _ {n} | \alpha_ {n} <   + \infty .
\]

  证明 A 是由 z 拓扑生成的交换幺 Banach 代数，且 A 的唯一极大理想是 \(x \in A\) 中无常数项的元素所成的集合。（注意 z 是拟幂零元。）

\begin{enumerate}
\def\labelenumi{\arabic{enumi})}
\setcounter{enumi}{4}
\item
  在代数 \(\mathbf{C}(\mathrm{X})\)（由 \(\mathbf{C}\) 上的有理分式组成）中，理想 \(\{0\}\) 是极大的，但其余维数不为 1。
\item
  设 \(\Delta\) 为 C 中的闭单位圆盘。证明 \(\mathcal{C}(\Delta)\) 中的可逆元群在 \(\mathcal{C}(\Delta)\) 中不稠密。
\item
  设 A 为交换幺 Banach 代数，\(\chi \in \mathsf{X}(\mathsf{A})\)。对每个 \(x \in \mathsf{A}\)，令 \(\|x\|_{\chi} = |\chi(x)| + \|x - \chi(x)\|\)。证明 \(x \mapsto \|x\|_{\chi}\) 是一个范数，\(\|xy\|_{\chi} \leqslant \|x\|_{\chi}\|y\|_{\chi}\)，\(\|e\|_{\chi} = 1\)，且若 \(\|e\| = 1\)，则 \(\|x\| \leqslant \|x\|_{\chi} \leqslant 3\|x\|\)。
\item
  设 A 为满足 \(\|1\| = 1\) 的交换幺 Banach 代数。设 N 为与给定范数等价、且使 A 仍是 Banach 代数、并使 1 的范数为 1 的范数全体所成的集合。
\end{enumerate}

\begin{enumerate}
\def\labelenumi{\alph{enumi})}
\tightlist
\item
  设 \(x_{0} \in A\) 满足 \(\varrho(x_{0}) < 1\)。对每个 \(x \in A\)，令：
\end{enumerate}

\[
\left\| x \right\| ^ {\prime} = \inf \sum_ {n} \left\| a _ {n} \right\|
\]

其中 \(x\) 取形如 \(a_0 + a_1 x_0 + \cdots + a_n x_0^n\) 的表示（\(a_0, a_1, \ldots, a_n \in \mathrm{A}\)）。证明 \(x \mapsto \|x\|^{\prime}\) 属于 N。（由于 \(\varrho(x_0) < 1\)，存在 \(k\) 使对一切 \(\|x_0^n\| \leqslant k\) 有 \(n\)，从而 \(\|x\| \leqslant k \|x\|^{\prime}\)。）

\begin{enumerate}
\def\labelenumi{\alph{enumi})}
\setcounter{enumi}{1}
\item
  证明 \(\|x_{0}\|^{\prime}\leqslant1\)。
\item
  由 a) 与 b) 推出，对每个 \(x \in A\)，有
\end{enumerate}

\[
\varrho (x) = \inf _ {n \in \mathrm{N}} n (x).
\]

\begin{enumerate}
\def\labelenumi{\alph{enumi})}
\setcounter{enumi}{3}
\tightlist
\item
  若 N 中唯一的元素就是给定的范数，则 \(A = C \cdot 1\)。（利用 c) 与习题 7。）
\end{enumerate}

\begin{enumerate}
\def\labelenumi{\arabic{enumi})}
\setcounter{enumi}{8}
\item
  设 X 为局部紧空间，A 为 Banach 代数 \(\mathcal{C}_{0}(X)\)。设 I 为 A 的一个理想，使得对每个 \(x \in X\)，存在 \(f \in I\) 满足 \(f(x) \neq 0\)。证明 I 包含 X 中具紧支撑的函数所成的理想。
\item
  设 A 为交换 Banach 代数，D 为 A 到 A 中的一个连续导子，即 A 到 A 中的连续线性映射，使得对 \(\mathrm{D}(ab)=(\mathrm{Da})b+a(\mathrm{Db})\) 有 \(a,b\in A\)。
\end{enumerate}

\begin{enumerate}
\def\labelenumi{\alph{enumi})}
\tightlist
\item
  对每个 \(\chi \in \mathsf{X}'(\mathbf{A})\) 与每个 \(\lambda \in \mathbf{C}\)，级数
\end{enumerate}

\[
\varphi_ {\lambda} (a) = \sum_ {n = 0} ^ {\infty} \frac {\lambda^ {n}}{n !} \chi (\mathrm{D} ^ {n} (a))
\]

收敛；\(\lambda \mapsto \varphi_{\lambda}(a)\) 是整函数，而 \(a \mapsto \varphi_{\lambda}(a)\) 是 A 的一个特征标。

\begin{enumerate}
\def\labelenumi{\alph{enumi})}
\setcounter{enumi}{1}
\item
  数 \(\varphi_{\lambda}(a)\) 与 \(\lambda\) 无关。（注意由 \(|\varphi_{\lambda}(a)| \leqslant \|a\|\)) 有 \(a\)。）由此得 \(\chi(\mathrm{D}(a)) = 0\)。
\item
  D 的像含于 A 的根之中。由此给出 I，第 162 页习题 31 的断言 f) 的一个新证明。
\end{enumerate}

\begin{enumerate}
\def\labelenumi{\arabic{enumi})}
\setcounter{enumi}{10}
\tightlist
\item
  设 B 为 {[}0,1{]} 上无限次可微的 C 值函数所成的代数。
\end{enumerate}

\begin{enumerate}
\def\labelenumi{\alph{enumi})}
\tightlist
\item
  设 A 为 B 的一个子代数，赋以一个使 A 成为 Banach 代数的范数。证明存在序列 \((m_{n})_{n\geqslant0}\)（取值于 \(R_{+}\)），使得对一切 \(x\in A\)，有
\end{enumerate}

\[
\sup _ {0 \leqslant t \leqslant 1} | x ^ {(n)} (t) | = \mathrm{O} (m _ {n}).
\]

（设 \(\mathrm{B}_n\) 为 \(n\) 次连续可微函数所成的 Banach 代数，定义在 {[}0,1{]} 上。由 I，第 40 页的命题 9，典范单射 \(\mathrm{A} \to \mathrm{B}_n\) 连续；设 \(\mathrm{M}_n\) 为其范数；若 \(x \in \mathrm{A}\)，则 \(\sup_{0 \leq t \leq 1} |x^{(n)}(t)| \leqslant n!\mathrm{M}_n \| x\|\)。）

\begin{enumerate}
\def\labelenumi{\alph{enumi})}
\setcounter{enumi}{1}
\tightlist
\item
  B 上不存在使它成为 Banach 代数的范数。
\end{enumerate}

\begin{enumerate}
\def\labelenumi{\arabic{enumi})}
\setcounter{enumi}{11}
\tightlist
\item
  设 A 为幺 Banach 代数，\((x_{n})\) 为 A 中趋于元素 \(x \in A\) 的可逆元序列。若序列 \((\varrho(x_{n}^{-1}))\) 有界，且对一切 n 有 \(x_{n}x = xx_{n}\)，则 x 可逆。（由 I，第 38 页命题 8 的推论，\(\varrho(1 - x_{n}^{-1}x) \leqslant \varrho(x_{n}^{-1})\varrho(x_{n} - x)\)，故对充分大的 n，\(x_{n}^{-1}x\) 可逆。）
\end{enumerate}

\begin{enumerate}
\def\labelenumi{¶ \arabic{enumi})}
\setcounter{enumi}{12}
\tightlist
\item
  设 A 为 Banach 代数 \(\ell^{2}(\mathbf{N})\)（习题 3）。设 \(\mathrm{A}_{0}\) 为由具有限支撑的序列组成的子代数。设 B 为 \(\mathrm{A}_{0}\) 与 \(\mathbf{C}\) 的直和，赋以乘法 \((f,\alpha)(g,\beta)=(fg,0)\)（对 \(f,g\in\mathrm{A}_{0},\alpha,\beta\in\mathbf{C}\)）以及范数
\end{enumerate}

\[
\left\| (f, \alpha) \right\| = \sup \left(\left\| f \right\|, \left| \alpha - \sum_ {n} f (n) \right|\right).
\]

\begin{enumerate}
\def\labelenumi{\alph{enumi})}
\item
  完备化代数 \(\widehat{\mathbf{B}}\) 是一个交换 Banach 代数，其根 R 为 \(\mathbf{C} \cdot (0,1)\)，且 \(\widehat{\mathbf{B}}/\mathbf{R}\) 等距同构于 A。
\item
  设 \(\pi\colon\widehat{\mathrm{B}}\to\widehat{\mathrm{B}}/\mathrm{R}\) 为典范态射。不存在 \(\mathrm{B}_{1}\)（\(\widehat{\mathrm{B}}\) 的、与 R 互补的子代数）使 \(\pi|\mathrm{B}_{1}\colon\mathrm{B}_{1}\to\widehat{\mathrm{B}}/\mathrm{R}=\mathrm{A}\) 是同胚。（设 \(\mathrm{B}_{1}\) 是这样一个子代数。设 \(u_{n}\in\mathrm{A}\) 满足 \(u_{n}(n)=1\)、\(u_{n}(m)=0\)（当 \(m\neq n\)）。设 \(e_{n}\in\mathrm{B}_{1}\) 满足 \(\pi(e_{n})=u_{n}\)。证明 \(e_{n}\in\mathrm{A}_{0}\)，且 \(\sum_{n=1}^{\infty}\frac{1}{n}u_{n}\) 在 A 中收敛，而 \(\sum_{n=1}^{\infty}\frac{1}{n}e_{n}\) 在 \(\widehat{\mathrm{B}}\) 中不收敛。）
\end{enumerate}

\begin{enumerate}
\def\labelenumi{\arabic{enumi})}
\setcounter{enumi}{13}
\tightlist
\item
  在 \(\mathbf{C}^2\) 中，设 \(\Omega\) 为由下式定义的紧集：
\end{enumerate}

\[
\frac {1}{2} \leqslant | z _ {1} | ^ {2} + | z _ {2} | ^ {2} \leqslant 1,
\]

并设 A 为 \(\mathcal{C}(\Omega)\) 的由在 \(\mathring{\Omega}\) 中全纯的函数组成的子代数。设 U 为由 \(|z_1|^2 + |z_2|^2 < 1\) 定义的开集。根据 Hartogs 定理（参见 J.-P. DEMAILLY, Complex analytic and differential geometry, Ch. I, th. 3.28），对每个 \(f \in \mathrm{A}\)，存在函数 \(\widetilde{f} \in \mathcal{C}(\overline{\mathrm{U}})\)

，它是唯一的，在 U 中全纯，且在 \(\Omega\) 上与 f 相合。证明 A 到 \(\mathcal{C}(\Omega)\) 中的典范单射是一个等距态射 h，其像是 \(\mathcal{C}(\Omega)\) 的一个全子代数，但 \(\mathsf{X}(h)\) 不是满射的。

\begin{enumerate}
\def\labelenumi{\arabic{enumi})}
\setcounter{enumi}{14}
\item
  设 A 与 B 为交换幺 Banach 代数，\(\varphi\) 为 A 到 B 的保单位元态射，\((x_{\lambda})_{\lambda \in \Lambda}\) 为 A 中一族元素，使得由诸 \(x_{\lambda}\) 生成的 A 的闭全子代数等于 A。为使 \(\mathsf{X}(\varphi)\) 是满射的，充要条件是 \(\mathrm{Sp}_{\mathrm{A}}^{\Lambda}((x_{\lambda})) = \mathrm{Sp}_{\mathrm{B}}^{\Lambda}((\varphi (x_{\lambda})))\)。（利用 I，第 41 页第 6 节的图 (1)，其中右端的箭头由 I，第 43 页的命题 12 知是双射。）
\item
  设 A 为交换 Banach 代数。下列条件等价：
\end{enumerate}

\begin{enumerate}
\def\labelenumi{(\roman{enumi})}
\item
  A 无根，且 \(\mathcal{G}(\mathrm{A})\) 在 \(\mathcal{C}_0(\mathsf{X}(\mathrm{A}))\) 中是闭的；
\item
  \(\mathcal{G}\) 是 A 到 \(\mathcal{G}(\mathrm{A})\) 上的一个同胚；
\item
  存在常数 c \textgreater{} 0，使得对一切 \(\|x\|^{2} \leqslant c\|x^{2}\|\) 有 \(x \in A\)。
\end{enumerate}

\begin{enumerate}
\def\labelenumi{\arabic{enumi})}
\setcounter{enumi}{16}
\item
  设 A 为交换 Banach 代数，I 为 A 的一个闭理想。用 h 表示 I 到 A 中的典范单射。空间 \(\mathsf{X}'(\mathsf{I})\) 典范地等同于 \(\mathsf{X}'(\mathsf{A})\) 关于由 \(\mathsf{X}'(h)\) 所定义的等价关系的商空间。（利用 I，第 153 页习题 4 的 a)。）
\item
  设 A 为交换幺 Banach 代数，\(x\) 为 A 的一个元素。假设由 \(x\) 生成的 A 的闭全子代数等于 A，于是映射 \(\chi \mapsto \chi(x)\) 使 \(\mathsf{X}(\mathsf{A})\) 与 \(\mathrm{Sp}_{\mathsf{A}}(x)\) 得以等同。设 \(\mathcal{H}\) 为函数 \(\mathcal{G}y\)（在 \(\mathrm{Sp}_{\mathsf{A}}(x)\) 上，其中 \(y\) 取遍 A）所成的集合。则 \(\check{\mathrm{S}}_{\mathcal{H}}(\mathrm{Sp}_{\mathsf{A}}(x))\)（INT，IV，§7，第 4 节）是 \(\mathrm{Sp}_{\mathsf{A}}(x)\) 相对于 C 的边界。（为证明 \(\check{\mathrm{S}}_{\mathcal{H}}(\mathrm{Sp}_{\mathsf{A}}(x))\) 含于这个边界，利用最大模原理。反之，设 \(z_0\) 为这个边界上的一点；为证明 \(z_0 \in \check{\mathrm{S}}_{\mathcal{H}}(\mathrm{Sp}_{\mathsf{A}}(x))\)，考虑 \((x - z_1)^{-1}\)，其中 \(z_1\) 充分接近于 \(z_0\)（在 \(\mathbf{C} = \mathrm{Sp}_{\mathsf{A}}(x)\) 中）。）
\item
  设 A 为交换幺 Banach 代数，B 为 A 的一个闭幺子代数，T 为从 X(A) 到 X(B) 中的映射 \(\chi \mapsto \chi|B\)，并设 R 为由 T 在 X(A) 上定义的等价关系。则 T 通过过渡到商，诱导出 X(A)/R 到 T(X(A)) 上的一个同胚，且对每个 \(x \in B\)，函数 \(f = \mathcal{G}_{\mathrm{B}}(x)\) 使 \(|f|\) 在 T(X(A)) 上达到其上确界。
\item
  设 \(\Delta\) 为圆盘 \(|z| \leqslant 1\)（在 \(\mathbf{C}\) 中），A 为由 \(f \in \mathcal{C}(\Delta)\) 中在 \(\Delta\) 的内部全纯的 f 所成的集合，而 \(\mathrm{A}_1\) 为由 A 与函数 \(\mathcal{C}(\Delta)\) 生成的 \(z \mapsto |z|\) 的 Banach 子代数。证明 \(\mathrm{X}(\mathrm{A}_1)\) 同胚于由 \((x_{1}, x_{2}, x_{3}) \in \mathbf{R}^{3}\)（满足 \(x_{1}^{2} + x_{2}^{2} \leqslant x_{3}^{2}\) 与 \(0 \leqslant x_{3} \leqslant 1\)）所成的集合。
\item
  \begin{enumerate}
  \def\labelenumii{\alph{enumii})}
  \tightlist
  \item
    设 \((\mathrm{X}_{\lambda})_{\lambda\in\Lambda}\) 为一族不定元。对每个形式级数 \(f\in\mathbf{C}[[(\mathrm{X}_{\lambda})]]_{\lambda\in\Lambda}\)，用 \(\|f\|\) 表示其系数的绝对值之和。设 A 为使 \(f\in\mathbf{C}[[(\mathrm{X}_{\lambda})]]_{\lambda\in\Lambda}\) 满足 \(\|f\|<+\infty\) 的那些 f 所成的集合。赋以通常的加法与乘法，A 是一个无根的交换幺 Banach 代数。
  \end{enumerate}
\end{enumerate}

\begin{enumerate}
\def\labelenumi{\alph{enumi})}
\setcounter{enumi}{1}
\tightlist
\item
  对每个交换幺 Banach 代数 B，存在一个 a) 中那种类型的代数 A，以及 A 到 B 上的一个连续保单位元态射。
\end{enumerate}

\begin{enumerate}
\def\labelenumi{\arabic{enumi})}
\setcounter{enumi}{21}
\tightlist
\item
  设 \(\Omega\) 为 \(\mathbf{C}^n\) 的一个紧子集。映射
\end{enumerate}

\[
(z _ {1}, \ldots , z _ {n}) \longmapsto (z _ {1}, \ldots , z _ {n}, \overline {{z}} _ {1}, \ldots , \overline {{z}} _ {n})
\]

是从 \(\Omega\) 到 \(\mathbf{C}^{2n}\) 中的一个紧的多项式凸子集上的同胚。（设 \(A = \mathcal{C}(\Omega)\)，\(z_i\) 为 \(\mathbf{C}^n\) 上的坐标函数。则 \((z_i|\Omega, \overline{z}_i|\Omega)_{1 \leqslant i \leqslant n}\) 是 A 的一个拓扑生成组，而所考虑的映射就是由这个拓扑生成组定义的、从 \(X(A) = \Omega\) 到 \(\mathbf{C}^{2n}\) 中的映射。）

\begin{enumerate}
\def\labelenumi{\arabic{enumi})}
\setcounter{enumi}{22}
\tightlist
\item
  用 \(\Omega\) 表示 \((z_1, z_2) \in \mathbf{C}^2\) 中满足 \(|z_1| \leqslant 1\)、\(|z_2| \leqslant 1\) 的元素所成的集合。设 \(\alpha \in [0,1]\)，并用 \(\Omega_\alpha\) 表示 \((z_1, z_2) \in \mathbf{C}^2\) 中满足
\end{enumerate}

\[
| z _ {1} | \leqslant 1, \quad | z _ {2} | \leqslant (1 - \alpha) | z _ {1} | + \alpha .
\]

\begin{enumerate}
\def\labelenumi{\alph{enumi})}
\item
  \(\Omega\) 与 \(\Omega_{\alpha}\) 在 \(\mathbf{C}^2\) 中的余集是连通的。
\item
  \(\Omega\) 是 \(\Omega_{\alpha}\) 的多项式凸包。（若 \((z_1^0, z_2^0) \in \Omega\) 且 \(\mathrm{P} \in \mathbf{C}[z_1, z_2]\)，则存在 \(z \in \mathbf{C}\) 使 \(|z| = 1\) 且 \(|\mathrm{P}(z_1^0, z_2^0)| \leqslant |\mathrm{P}(z, z_2^0)|\)；而 \(|\mathrm{P}(z, z_2^0)| \leqslant \sup_{(z_1, z_2) \in \Omega_{\alpha}} |\mathrm{P}(z_1, z_2)|\)。）
\end{enumerate}

\begin{enumerate}
\def\labelenumi{\arabic{enumi})}
\setcounter{enumi}{23}
\item
  设 \(n \geqslant 1\)。存在 \(\mathbf{C}^{n}\) 的一个闭凸子集，它不是多项式凸的。
\item
  设 A 为一个完备的交换幺局部 m-凸代数（参见 I，第 165 页，习题 37）。
\end{enumerate}

\begin{enumerate}
\def\labelenumi{\alph{enumi})}
\item
  A 的元素 x 可逆，充要条件是 \(\chi(x) \neq 0\) 对 A 的每个连续特征标 \(\chi\) 成立。A 的根是 A 的诸连续特征标的核的交，故是闭的。
\item
  在 I，第 165 页习题 37 的例 e) 中，设 I 为由使对每个充分大的整数 n \(f(n) = 0\)\textgreater{ 0 都有 } 的 f 所成的理想。则每个包含 I 的极大理想都是稠密的且具有无穷余维数。
\end{enumerate}

\begin{enumerate}
\def\labelenumi{\arabic{enumi})}
\setcounter{enumi}{25}
\tightlist
\item
  设 X 为紧拓扑空间，E 为复 Banach 空间，其范数记作 \(\|\cdot\|_{E}\)。用 \(\mathcal{C}(X,E)\) 表示从 X 到 E 中的连续函数所成的复向量空间。
\end{enumerate}

\begin{enumerate}
\def\labelenumi{\alph{enumi})}
\tightlist
\item
  赋以范数
\end{enumerate}

\[
\| f \| = \sup _ {x \in \mathrm{X}} \| f (x) \| _ {\mathrm{E}},
\]

空间 \(\mathcal{C}(\mathrm{X},\mathrm{E})\) 是一个 Banach 空间。

\begin{enumerate}
\def\labelenumi{\alph{enumi})}
\setcounter{enumi}{1}
\tightlist
\item
  设 B 为 \(\mathcal{C}(\mathrm{X},\mathrm{E})\) 的由函数 \(x\mapsto f(x)y\)（其中 \(f\in \mathcal{C}(\mathrm{X})\)、\(y\in \mathrm{E}\)）生成的向量子空间。证明 B 在 \(\mathcal{C}(\mathrm{X},\mathrm{E})\) 中稠密。
\end{enumerate}

在以下内容中，我们考虑一个交换 Banach 代数 A。

\begin{enumerate}
\def\labelenumi{\alph{enumi})}
\setcounter{enumi}{2}
\tightlist
\item
  赋以乘法 \((fg)(x) = f(x)g(x)\)（对一切 \(x\in \mathrm{X}\)），Banach 空间 \(\mathcal{C}(\mathrm{X},\mathrm{A})\) 是一个交换 Banach 代数。
\end{enumerate}

\begin{enumerate}
\def\labelenumi{\alph{enumi})}
\setcounter{enumi}{3}
\tightlist
\item
  每个特征标 \(\chi\in\mathsf{X}(\mathcal{C}(\mathrm{X},\mathrm{A}))\) 都具有形式 \(f\mapsto\widetilde{\chi}(f(x))\)，其中 \(\widetilde{\chi}\in\mathsf{X}(\mathrm{A})\) 且 \(x\in\mathrm{X}\)。（先处理 A 幺且 \(\|1\|=1\) 的情形；考虑 \(\chi\) 在 \(\mathcal{C}(\mathrm{X})\cdot1\) 上的限制，以及它在同态 \(\mathrm{A}\to\mathcal{C}(\mathrm{X},\mathrm{A})\) 的像上的限制，该同态把 \(a\) 映为常值函数 \(x\mapsto a\)。）
\end{enumerate}

\begin{enumerate}
\def\labelenumi{\alph{enumi})}
\setcounter{enumi}{4}
\tightlist
\item
  从 \(\mathsf{X}(\mathrm{A})\times \mathrm{X}\) 到 \(\mathsf{X}(\mathcal{C}(\mathrm{X},\mathrm{A}))\) 中的、把 \((\widetilde{\chi},x)\) 对应到特征标 \(f\mapsto \widetilde{\chi} (f(x))\) 的映射是一个同胚。
\end{enumerate}

\begin{enumerate}
\def\labelenumi{\arabic{enumi})}
\setcounter{enumi}{26}
\tightlist
\item
  设 A 为 C 上的交换幺 Banach 代数。X(A) 的闭子集 F 称为 A 的边界，如果有
\end{enumerate}

\[
\sup _ {\chi \in \mathsf {X} (\mathrm{A})} | \chi (x) | = \sup _ {\chi \in \mathrm{F}} | \chi (x) |
\]

对所有 \(x \in A\) 成立。

\begin{enumerate}
\def\labelenumi{\alph{enumi})}
\item
  将 A 的边界所成的集合按包含关系排序。存在一个极小边界。
\item
  设 S 为 A 的一个极小边界。A 的每个边界都包含 S，因而 S 是 A 的所有边界之交。（设 F 为 A 的一个边界；证明对每个 \(x \in S\)，x 的每个邻域 U 都与 F 相交；为此，利用 S − U 不是 A 的边界这一事实。）
\end{enumerate}

称 S 为 A 的 Shilov 边界，并记作 \(\partial A\)。

\begin{enumerate}
\def\labelenumi{\alph{enumi})}
\setcounter{enumi}{2}
\item
  若 A 的 Gelfand 变换的像在 \(\mathcal{C}(\mathrm{X}(\mathrm{A}))\) 中稠密，则 A 的 Shilov 边界等于 \(\mathrm{X}(\mathrm{A})\)。
\item
  设 A 为单位圆盘 \(\Delta\)（\(z \in C\) 满足 \(|z| \leqslant 1\)）上的、在其内部 \(\Delta\) 全纯的连续函数所成的代数（I，第 20 页的例 9）。则 \(\mathsf{X}(\mathsf{A})\) 典范地等同于 \(\Delta\)（I，第 193 页的习题 6），且 A 的 Shilov 边界等同于单位圆周 U。
\end{enumerate}

\begin{enumerate}
\def\labelenumi{\arabic{enumi})}
\setcounter{enumi}{27}
\item
  设 X 为紧拓扑空间，n 为整数 \(\geqslant 1\)。设 I 为 Banach 代数 \(\mathcal{C}(\mathrm{X},\mathrm{M}_{n}(\mathbf{C}))\) 的一个双边理想。证明存在唯一的闭子集 S ⊂ X，使得 I 恰为 \(f \in \mathcal{C}(\mathrm{X}, \mathrm{M}_{n}(\mathbf{C}))\) 中在 S 上限制为零的 f 所成的集合。
\item
  * 设 X 为紧实流形。给 \(\mathcal{C}^{\infty}(\mathrm{X})\) 赋以函数及其各阶导数一致收敛的拓扑；它是一个 Fréchet 空间。设 A 为 \(\mathcal{C}(\mathrm{X})\) 的一个子代数。假设 A 赋有一个使它成为 Banach 代数的范数 \(f \mapsto \| f\|_{\mathrm{A}}\)，且 A 到 \(\mathcal{C}(\mathrm{X})\) 中的典范单射连续。再假设 \(\mathcal{C}^{\infty}(\mathrm{X})\) 含于 A 中且在 A 中稠密。
\end{enumerate}

\begin{enumerate}
\def\labelenumi{\alph{enumi})}
\tightlist
\item
  存在常数 \(C \geqslant 0\) 与整数 \(k \in N\)，使得对一切 \(f \in \mathcal{C}^{\infty}(X)\)，有
\end{enumerate}

\[
\| f \| _ {\mathrm{A}} \leqslant \operatorname * {C} \sup _ {x \in \mathrm{X}} | f ^ {(k)} (x) |.
\]

\begin{enumerate}
\def\labelenumi{\alph{enumi})}
\setcounter{enumi}{1}
\item
  子代数 A 在 \(\mathcal{C}(\mathrm{X})\) 中是全的。（设 \(f \in \mathrm{A}\) 在 X 上不取零值，并设 \(\delta = \inf_{x \in \mathrm{X}} |f(x)|\)；选取 \(g \in \mathcal{C}^{\infty}(\mathrm{X})\) 使 \(\|f - g\|_{\mathrm{A}} < \delta/3\)，然后在 \(f^{-1} = g^{-1} \sum_{n \in \mathbf{N}} ((g - f)/g)^{n}\) 中写出 \(\mathcal{C}(\mathrm{X})\)，并验证该级数在 A 中收敛。）
\item
  由此给出 Wiener 定理（I，第 38 页的例）的一个新证明。
\end{enumerate}

\exercisesection{全纯函数演算}

在以下习题中，所考虑的代数除特别声明外均以 C 为基域。

\begin{enumerate}
\def\labelenumi{\arabic{enumi})}
\tightlist
\item
  设 \(n \geqslant 1\) 为整数，且 \(A = \mathcal{L}(\mathbf{C}^{n})\)。设 \(\alpha \in C\)，并设 \(x \in A\) 为一个元素，它相对于 \(C^{n}\) 的典范基的矩阵是 n 阶、以 \(\alpha\) 为特征值的 Jordan 矩阵。x 的谱化为 \(\alpha\)。
\end{enumerate}

\begin{enumerate}
\def\labelenumi{\alph{enumi})}
\tightlist
\item
  对 \(f \in \mathcal{O}(\{\alpha\})\)，计算相对于 \(\mathbf{C}^n\) 的典范基、表示 \(f(x)\) 的矩阵。
\end{enumerate}

  设 \(m \geqslant 1\) 为整数。设 \(B_{m}\) 为 \(C^{m}\) 类函数在 \(\alpha\) 的一个邻域中的芽所成的幺 Banach 代数，赋以 \(C^{m}\) 一致收敛拓扑（VAR，R2，§12）。

\begin{enumerate}
\def\labelenumi{\alph{enumi})}
\setcounter{enumi}{1}
\item
  存在一个从 \(B_{m-1}\) 到 A 的连续幺代数态射 \(\varphi(z)=x\)，其中 \(z\in B_{m-1}\) 是 C 的恒等函数在 \(\alpha\) 处的 m−1 阶芽。
\item
  不存在连续幺代数态射 \(\varphi\)，从 \(B_m\) 到 A，使 \(\varphi(z) = x\)，其中 \(z \in B_m\) 是 \(\mathbf{C}\) 的恒等函数的芽。
\end{enumerate}

\begin{enumerate}
\def\labelenumi{\arabic{enumi})}
\setcounter{enumi}{1}
\tightlist
\item
  设 A 为幺 Banach 代数，G 为 A 的可逆元群。
\end{enumerate}

\begin{enumerate}
\def\labelenumi{\alph{enumi})}
\item
  为使元素 \(x \in A\) 具有形式 \(\exp(y)\)（其中 \(y \in A\)），充要条件是它含于 G 的一个连通交换子群之中。
\item
  为使元素 \(x \in \mathbf{A}\) 具有形式 \(\exp(y)\)（其中 \(y \in \mathbf{A}\)），只需 0 属于 \(\mathbf{C} = \mathrm{Sp}_{\mathrm{A}}(x)\) 的无界连通分支。
\end{enumerate}

\begin{enumerate}
\def\labelenumi{\arabic{enumi})}
\setcounter{enumi}{2}
\tightlist
\item
  设 A 为幺 Banach 代数，G 为 A 的可逆元群，\(G_{1}\) 为 G 的单位元连通分支。
\end{enumerate}

\begin{enumerate}
\def\labelenumi{\alph{enumi})}
\item
  设 \(x \in \mathbf{A}\)，\(n\) 为整数 \(\geqslant 0\)，使 \(x^n = e\)。证明 \(x \in \mathbf{G}_1\)。（\(x\) 的谱有限。使 \(\lambda \in \mathbf{C}\) 满足 \(y(\lambda) = \lambda x + e - \lambda e\) 为可逆者所成的集合，其余集有限，从而连通。而 \(y(0) = e\)，\(y(1) = x\)。
\item
  假设 A 交换。G/G₁ 中每个异于单位元的元素都有无穷阶。（若 \(x \in G\) 且 \(x^n \in G_1\)，则由 I，第 78 页第 10 节，有 \(x^n = \exp(y)\)（其中 \(y \in A\)），从而 \(x \exp(-y/n)^n = e\)。应用 a)。）
\end{enumerate}

\begin{enumerate}
\def\labelenumi{\arabic{enumi})}
\setcounter{enumi}{3}
\item
  存在一个交换 Banach 代数，非幺且非零，使得 \(X(A)\) 是紧的。
\item
  设 A 为幺 Banach 代数，\(a \in A\)。设 G 为 A 的可逆元群，\(G_1\) 为 G 中含单位元的连通分支。A 的指数谱是使 \(\lambda \in \mathbf{C}\) 不属于 \(x - \lambda e\) 的 \(G_1\) 所成的集合（参见 I，第 79 页的命题 12）。
\end{enumerate}

\begin{enumerate}
\def\labelenumi{\alph{enumi})}
\item
  a 的指数谱是闭的；它包含 a 的谱。
\item
  a 的指数谱的边界含于 a 的奇异谱（I，第 158 页的习题 16）。
\item
  a 的指数谱是 a 的谱与 \(\mathbf{C}=\mathrm{Sp}(a)\) 的某些有界连通分支的并。
\end{enumerate}

\begin{enumerate}
\def\labelenumi{\alph{enumi})}
\setcounter{enumi}{3}
\tightlist
\item
  设 \(f\in \mathbf{C}[\mathrm{X}]\) 为多项式。\(f\) 在 \(a\) 的指数谱上的像包含 \(f(a)\) 的指数谱。一般而言，此包含是严格的（利用 Fredholm 理论，参见 III，待出版）。
\end{enumerate}

\begin{enumerate}
\def\labelenumi{\arabic{enumi})}
\setcounter{enumi}{5}
\tightlist
\item
  设 A 为交换幺 Banach 代数。
\end{enumerate}

\begin{enumerate}
\def\labelenumi{\alph{enumi})}
\item
  若 j 是 A 的一个幂等元，则 \(\exp(2i\pi j)=1\) 且 \(\exp(i\pi j)=1-2j\)。
\item
  设 \(p \in A\) 使 \(\exp(p) = 1\)。则 \(\mathrm{Sp}_{\mathrm{A}}(p)\) 由形如 \(2in\pi\)（其中 \(n \in Z\)）的有限多个点组成。
\item
  对 \(\lambda \notin \mathrm{Sp}_{\mathrm{A}}(p)\)，我们有：
\end{enumerate}

\[
\int_ {0} ^ {1} \exp (t (p - \lambda 1)) d t = (1 - \exp (- \lambda)) (\lambda 1 - p) ^ {- 1}.
\]

\begin{enumerate}
\def\labelenumi{\alph{enumi})}
\setcounter{enumi}{3}
\tightlist
\item
  由此推出 \(\mathrm{Sp}_{\mathbf{A}}(p)\) 的每一点都是预解式 \(\mathrm{R}(p,\lambda)\) 的单极点，进而存在整数 \(m \geqslant 0\)、整数 \(n_{k} \in Z\) 以及两两正交的幂等元 \(j_{k}\)（对 \(1 \leqslant k \leqslant m\)），使得
\end{enumerate}

\[
p = 2 i \pi \sum_ {k} n _ {k} j _ {k}.
\]

\begin{enumerate}
\def\labelenumi{¶ \arabic{enumi})}
\setcounter{enumi}{6}
\tightlist
\item
  a) 设 B 为复 Banach 空间，j 为 B 中满足 \(\|j\| = 1\) 的一点。对每个 \(x \in B\)，极限
\end{enumerate}

\[
\lim_{\substack{\alpha \to 0\\ \alpha >0}}\frac{1}{\alpha} (\| j + \alpha x\| -1)
\]

存在。令 \(\varphi(x)\) 为这个极限。

\begin{enumerate}
\def\labelenumi{\alph{enumi})}
\setcounter{enumi}{1}
\tightlist
\item
  设
\end{enumerate}

\[
\psi (x) = \sup_{\substack{\varepsilon \in \mathbf{C}\\ |\varepsilon | = 1}}\varphi (\varepsilon x).
\]

函数 \(\psi\) 是一个被 B 的范数所控制的半范数。我们有

\[
\psi (x) = \sup _ {f} | f (x) |,
\]

其中 \(f\) 取遍 B 的对偶中满足 \(\|f\| = f(j) = 1\) 的元素所成的集合。

\begin{enumerate}
\def\labelenumi{\alph{enumi})}
\setcounter{enumi}{2}
\tightlist
\item
  现在假设 B 是一个 Banach 代数，且 \(j = 1\) 是 B 的单位元。我们有：
\end{enumerate}

\[
\varphi (x) = \lim_{\substack{\alpha \to 0\\ \alpha >0}}\frac{1}{\alpha}\log (\| \exp (\alpha x)\|) = \sup_{\alpha >0}\frac{1}{\alpha}\log (\| \exp (\alpha x)\|).
\]

（注意 \(\log(\|\exp((\alpha+\alpha')x)\|) \leqslant \log(\|\exp(\alpha x)\|) + \log(\|\exp(\alpha' x)\|).\)。）

\begin{enumerate}
\def\labelenumi{\alph{enumi})}
\setcounter{enumi}{3}
\tightlist
\item
  利用 c)，证明 \(\|\exp(\lambda x)\| \leqslant \exp(|\lambda|\psi(x))\)。在由 \(\lambda \mapsto \lambda^{-2} \exp(\lambda x)\) 所确定的圆周上积分 \(|\lambda| = \varrho\)，推出
\end{enumerate}

\[
\| x \| \leqslant \varrho^ {- 1} \exp (\varrho \psi (x)),
\]

从而令 \(\varrho = \psi(x)^{-1}\)，得 \(\|x\| \leqslant e\psi(x)\)，其中 \(e = \exp(1)\)。

\begin{enumerate}
\def\labelenumi{\alph{enumi})}
\setcounter{enumi}{4}
\tightlist
\item
  对 B 的每个非零元素 \(x\)，存在线性形式 \(f \in \mathrm{B}'\)，使 \(f(1) = \| f \| = 1\) 且 \(f(x) \neq 0\)。
\end{enumerate}

\begin{enumerate}
\def\labelenumi{\arabic{enumi})}
\setcounter{enumi}{7}
\item
  设 A 为幺 Banach 代数，\(x \in A\)，B 为由 x 生成的 A 的子代数。为使 \(\mathrm{Sp}_{\mathrm{A}}(x)\) 是一个有限集且其所有元素都是 x 的预解式的极点，充要条件是 B 是有限维的。（为证条件必要：若 \(\mathrm{R}(x, \lambda) = \sum_{i=1}^{p} \mathrm{F}_{i}(\lambda) a_{i}\)（其中 \(a_{i} \in A\)，\(F_{i}\) 为纯量有理函数），则 \(\mathrm{R}(x, \lambda)\) 在 \(\lambda = \infty\) 处的展开式表明诸 \(x^{n}\) 是诸 \(a_{i}\) 的线性组合。）
\item
  设 A 为幺 Banach 代数，\(x \in A\)，U 为 \(\mathrm{Sp}_{\mathrm{A}}(x)\) 的一个只有有限多个连通分支 \(U_{i}\)（\(i \in I\)）的开邻域，且 \(f \in \mathcal{O}(\mathrm{U})\)。为使 \(f(x) = 0\)，充要条件是下列条件被满足：
\end{enumerate}

\begin{enumerate}
\def\labelenumi{(\roman{enumi})}
\item
  对每个使 \(\mathrm{Sp}_{\mathrm{A}}(x)\cap\mathrm{U}_{i}\) 为无限集或含有一个不是预解式极点的点的 i，有 \(f|U_{i}=0\)；
\item
  预解式的每个 \(p\) 阶极点都是 \(\geqslant p\) 的阶 \(f\) 的零点。（关于 (ii) 的必要性，按习题 8 的方式论证。）
\end{enumerate}

\begin{enumerate}
\def\labelenumi{\arabic{enumi})}
\setcounter{enumi}{9}
\tightlist
\item
  设 \(\Delta\) 为 C 中的圆盘 \(|z| \leqslant 1\)。设 A 为由在 \(\Delta\) 上连续、在 \(\Delta\) 的内部全纯、且存在可和的复数序列 \((c_{n})\) 满足
\end{enumerate}

\[
f (z) = \sum_ {n \geqslant 0} c _ {n} z ^ {n}
\]

（对一切 \(z \in \Delta\)）的函数 f 所成的空间。这时记 \(\|f\| = \sum_{n}|c_{n}|\)。

\begin{enumerate}
\def\labelenumi{\alph{enumi})}
\item
  赋以乘法 \((f,g)\mapsto fg\)，空间 A 是一个幺 Banach 代数，且 z 是 A 的一个拓扑生成元。
\item
  空间 \(\mathsf{X}(\mathsf{A})\) 典范地等同于 \(\Delta\)；若 \(f \in A\) 且 g 在 \(f(\Delta)\) 的一个开邻域中全纯，则 \(g \circ f \in A\)。
\end{enumerate}

\begin{enumerate}
\def\labelenumi{\arabic{enumi})}
\setcounter{enumi}{10}
\item
  设 A 为交换幺 Banach 代数，\(x \in A\)。设 S 为 X(A) 的一个对 Jacobson 拓扑闭的子集，且 f 为在 \((\mathcal{G}x)(\mathrm{S})\) 的一个邻域中全纯的复值函数。存在 \(y \in A\) 使在 S 上 \(Gy = f \circ Gx\)。（设 I 为诸 \(\chi \in S\) 的核的交。在 A/I 中使用泛函演算。）
\item
  设 U 为 C 中的非空开集。证明 \(\mathcal{O}(\mathrm{U})\) 上不存在使 \(\mathcal{O}(\mathrm{U})\) 成为 Banach 代数的范数。（利用 I，第 29 页的定理 1，证明 \(\mathcal{O}(\mathrm{U})\) 中的函数都将有界。）
\item
  设 \(n\in \mathbf{N}\)
\end{enumerate}

\begin{enumerate}
\def\labelenumi{\alph{enumi})}
\tightlist
\item
  \(C^{n}\) 的子集 V 是多项式凸的，充要条件是：对每个 \(x \in C^{n} - V\)，存在整函数 \(f: C^{n} \to C\) 使得
\end{enumerate}

\[
| f (x) | > \sup _ {y \in V} | f (y) |.
\]

\begin{enumerate}
\def\labelenumi{\alph{enumi})}
\setcounter{enumi}{1}
\item
  设 \(K \subset R^{n}\) 为紧集。则 K 在 \(C^{n}\) 中是多项式凸的。
\item
  设 V 为 \(C^{n}\) 的一个多项式凸紧子集，\(f \in \mathcal{O}(V)\)。则 f 的图象在 \(C^{n+1}\) 中是多项式凸的。（利用 Oka–Weil 定理。）
\end{enumerate}

\begin{enumerate}
\def\labelenumi{\arabic{enumi})}
\setcounter{enumi}{13}
\item
  设 K 为 C 的紧子集，P（相应地，Q）为在 K 的邻域中全纯的多项式（相应地，有理）函数的芽所成集合在 \(\mathcal{O}(\mathrm{K})\) 中的闭包。证明 \(\mathrm{P} = \mathrm{Q}\) 的充要条件是 \(\mathbf{C} - \mathbf{K}\) 连通。
\item
  设 U 为模为 1 的复数所成的集合。设 \(K_{1}\)、\(K_{2}\) 为 \(C^{2}\) 中由下式给出的紧子集：
\end{enumerate}

\[
\mathrm{K} _ {1} = \mathbf {U} \times \{0 \}, \quad \mathrm{K} _ {2} = \{(z, \overline {{z}}) \mid z \in \mathbf {U} \}.
\]

\begin{enumerate}
\def\labelenumi{\alph{enumi})}
\item
  证明：在 \(K_{1}\) 的邻域中的、C 上复系数多项式函数的芽所成的集合在 \(\mathcal{O}(K_{1})\) 中不稠密。
\item
  证明：在 \(K_{2}\) 的邻域中的、C 上复系数多项式函数的芽所成的集合在 \(\mathcal{O}(K_{2})\) 中稠密。
\end{enumerate}

\begin{enumerate}
\def\labelenumi{\arabic{enumi})}
\setcounter{enumi}{15}
\tightlist
\item
  设 A 为幺 Banach 代数，\(x \in A\)，k 为整数 \(\geqslant 0\)。设 U 为模为 1 的复数所成的集合。
\end{enumerate}

\begin{enumerate}
\def\labelenumi{\alph{enumi})}
\tightlist
\item
  假设 \(\mathrm{Sp}_{\mathrm{A}}(x)\subset\mathbf{U}\)，且当 \(\|\mathrm{R}(x,\lambda)\|=\mathrm{O}((1-|\lambda|)^{k})\) 趋于 1 时 \(|\lambda|\)。证明当 n 趋于 \(\|x^{n}\|=\mathrm{O}(|n|^{k})\) 时 \(\pm\infty\)。（对每个 \(\delta>1\)，证明
\end{enumerate}

\[
2 \pi \| x ^ {n} \| \leqslant \mathrm{M} \frac {\delta^ {n}}{(\delta - 1) ^ {k}} + \mathrm{M} \frac {\delta^ {- n}}{(1 - \delta^ {- 1}) ^ {k}}
\]

其中 M 与 n 和 \(\delta\) 无关。当 n 趋于 \(\delta = n/(n - k)\) 时取 \(+\infty\)；当 n 趋于 \(\delta = n/(n + k)\) 时取 \(-\infty\)。

\begin{enumerate}
\def\labelenumi{\alph{enumi})}
\setcounter{enumi}{1}
\tightlist
\item
  假设当 n 趋于 \(\|x^{n}\|=\mathrm{O}(|n|^{k})\) 时 \(\pm\infty\)。证明 \(\operatorname{Sp}_{\mathbf{A}}(x)\subset\mathbf{U}\)，且当 \(\|\operatorname{R}(\lambda,x)\|=\operatorname{O}((1-|\lambda|)^{k+1})\) 趋于 1 时 \(|\lambda|\)。（我们有 \(\varrho(x)=\varrho(x^{-1})=1\)，从而 \(\operatorname{Sp}_{\mathbf{A}}(x)\subset\mathbf{U}\)。对 \(|\lambda|<1\)，
\end{enumerate}

\[
\mathrm{R} (\lambda , x) = - x ^ {- 1} \left(1 + \lambda x ^ {- 1} + \lambda^ {2} x ^ {- 2} + \lambda^ {3} x ^ {- 3} + \dots\right).
\]

由此导出 \(\|\mathrm{R}(\lambda,x)\|\) 的一个上界估计，并对 \(|\lambda|>1\) 的情形类似地论证。

\begin{enumerate}
\def\labelenumi{\arabic{enumi})}
\setcounter{enumi}{16}
\tightlist
\item
  设 \(X_{1}\)（相应地，\(X_{2}\)）为 \((z_{1}, z_{2}) \in \mathbf{C}^{2}\) 中满足下式的元素所成的集合：
\end{enumerate}

\[
\frac {1}{2} \leqslant | z _ {1} | ^ {2} + | z _ {2} | ^ {2} \leqslant 1, \quad \text { resp. } \quad | z _ {1} | ^ {2} + | z _ {2} | ^ {2} \leqslant 1.
\]

  设 \(X = X_{1} \cup \{(0,0)\} \subset \mathbf{C}^{2}\)，\(A = \mathcal{C}(X)\)，并设 \(a_{1}, a_{2}\) 为 \(C^{2}\) 上坐标函数在 X 上的限制。

\begin{enumerate}
\def\labelenumi{\alph{enumi})}
\item
  联合谱 \(\mathrm{Sp}_{\mathrm{A}}^{2}(a_{1},a_{2})\) 等于 X。
\item
  构造两个不同的、从 \(\mathcal{O}(\mathrm{X})\) 到 A 的连续保单位元态射，它们把 \(z_{1}\) 映为 \(a_{1}\)，把 \(z_{2}\) 映为 \(a_{2}\)。（利用下述定理：对每个在 \(f\) 的邻域中全纯的函数 \(\mathrm{X}_{1}\)，存在在 \(X_{2}\) 的邻域中全纯且在 \(X_{1}\) 的邻域中与 f 相合的函数 g（参见 I，第 168 页，习题 14）。
\end{enumerate}

\begin{enumerate}
\def\labelenumi{¶ \arabic{enumi})}
\setcounter{enumi}{17}
\tightlist
\item
  设 A 为交换幺 Banach 代数。
\end{enumerate}

\begin{enumerate}
\def\labelenumi{\alph{enumi})}
\tightlist
\item
  对 A 的元素的任何有限族 I，设 S(I) ⊂ C\(^{I}\) 为 I 的联合谱。若 I′ 是 I 的扩张，则典范投射 pr\(_{I,I'}\)：从 C\(^{I'}\) 到 C\(^{I}\) 上的投影，满足 pr\(_{I,I'}\)(S(I′)) = S(I)，从而得到 \(\varphi_{II'}\)：从 \(\mathcal{O}(S(I))\) 到 \(\mathcal{O}(S(I'))\) 中的单射态射。设 \(\mathcal{O}(X(A))\) 为诸 \(\mathcal{O}(S(I))\) 关于诸 \(\varphi_{II'}\) 的归纳极限（我们用 A × N 的有限子集来标记 A 的元素的有限族）。典范映射
\end{enumerate}

\[
\varphi_ {\mathrm{I}} \colon \mathcal {O} (\mathrm{S} (\mathrm{I})) \longrightarrow \mathcal {O} (\mathrm{X} (\mathrm{A}))
\]

都是单射的。

我们将给 \(\mathcal{O}(\mathrm{X(A)})\) 赋以诸 \(\mathcal{O}(\mathrm{S(I)})\) 的拓扑的归纳极限拓扑。对每个 \(a \in \mathrm{A}\)，\(\mathrm{S}(a) \subset \mathbf{C}\) 在 C 中某邻域内的坐标函数定义了 \(z_a\) 的一个元素 \(\mathcal{O}(\mathrm{X(A)})\)。

\begin{enumerate}
\def\labelenumi{\alph{enumi})}
\setcounter{enumi}{1}
\tightlist
\item
  连续满射 \(\mathsf{X}(\mathsf{A}) \to \mathsf{S}(\mathsf{I})\) 定义了连续态射 \(\mathcal{C}(\mathsf{S}(\mathsf{I})) \to \mathcal{C}(\mathsf{X}(\mathsf{A}))\)。态射 \(\mathcal{O}(\mathsf{S}(\mathsf{I})) \to \mathcal{C}(\mathsf{S}(\mathsf{I})) \to \mathcal{C}(\mathsf{X}(\mathsf{A}))\) 定义了一个连续态射：
\end{enumerate}

\[
\varphi \colon \mathcal {O} (\mathrm{X} (\mathrm{A})) \longrightarrow \mathcal {C} (\mathrm{X} (\mathrm{A})).
\]

这个态射一般不是单射的。

\begin{enumerate}
\def\labelenumi{\alph{enumi})}
\setcounter{enumi}{2}
\item
  存在唯一的从 \(\psi\) 到 A 的连续保单位元态射 \(\mathcal{O}(\mathsf{X}(\mathsf{A}))\)，使对一切 \(\psi(z_a) = a\) 有 \(a \in \mathsf{A}\)。（利用 I，第 51 页的定理 1。）它与 Gelfand 变换的复合是 \(\varphi\)。
\item
  设 E 为分离局部凸空间，U 为 E 的弱开子集，\(f: U \to C\) 为一个函数。若对每个 \(u \in U\)，存在 u 的一个弱开邻域 \(V_u\)、E 中余维数有限的闭向量子空间 \(F_u\)、以及 \(g_u\) 上的全纯函数 \(p_u(V_u)\)（其中 \(p_u\) 表示 E 到 \(E/F_u\) 上的典范映射），使得 \(f|V_u = g_u \circ p_u\)，则称 f 是弱全纯的。
\end{enumerate}

  证明：若 \(f\) 弱全纯，且 K 是 U 的一个弱紧子集，则存在 K 在 U 中的一个弱开邻域 \(V_K\)、E 中余维数有限的闭向量子空间 \(F_K\)、以及 \(g_K\) 上的全纯函数 \(p_K(V_K)\)，使得 \(f|V_K = g_K \circ p_K\)，其中 \(p_K\) 是 E 到 \(E/F_K\) 上的典范映射。

\begin{enumerate}
\def\labelenumi{\alph{enumi})}
\setcounter{enumi}{4}
\tightlist
\item
  证明 \(\mathcal{O}(\mathsf{X}(\mathsf{A}))\) 是诸 \(\mathcal{O}(\mathsf{S})\) 的归纳极限，其中 S 取遍 A 的对偶中 \(\mathsf{X}(\mathsf{A})\) 的弱开邻域所成的集合，而 \(\mathcal{O}(\mathsf{S})\) 表示 S 上弱全纯函数所成的代数，赋以紧收敛拓扑。
\end{enumerate}

\exercisesection{正则交换 Banach 代数}

在以下习题中，所考虑的代数均以 C 为基域。

\begin{enumerate}
\def\labelenumi{\arabic{enumi})}
\item
  设 A 为正则交换 Banach 代数，I 为 A 的一个闭理想。则 I 与 A/I 都是正则交换 Banach 代数。
\item
  设 A 为无根正则交换 Banach 代数，I 为 A 的理想，\(\chi\) 为 V(I) 的元素，J 为由使 Gx 的支撑紧且与 \(x \in A\) 不相交的 \(\{\chi\}\) 所成的集合。则 I + J 是 A 的理想 R 中使 R ⊃ I 且 V(R) = \(\{\chi\}\) 的最小者。
\item
  设 A 为交换 Banach 代数，A′ 为 A 的对偶 Banach 空间。对每个 \(x \in A\)，用 \(\kappa(x)\) 表示 A′ 中与 A 中乘以 x 的算子相转置的线性算子。对 \(x \in A\) 与 \(x' \in A'\)，记 \(x \star x' = \kappa(x)(x') \in A'\)。若 A′ 的一个向量子空间在所有 \(\kappa(x)\) 之下稳定，则称它是不变的。
\end{enumerate}

\begin{enumerate}
\def\labelenumi{\alph{enumi})}
\item
  为使 A′ 的元素 \(x'\) 与 A 的一个非零特征标成比例，充要条件是 \(Cx'\) 是不变的。
\item
  映射 \(V \mapsto V^{\circ}\) 是从 A 的闭理想集合到 \(A'\) 的弱闭不变向量子空间集合上的一个双射。
\item
  若 W 是 A′ 的不变向量子空间，用 \(\sigma(\mathrm{W})\) 表示属于 W 的特征标所成的集合。若 \(W_{1}\) 与 \(W_{2}\) 是弱闭不变子空间，且 W 是由 \(W_{1}\) 与 \(W_{2}\) 生成的弱闭向量子空间，则 \(\sigma(\mathrm{W}) = \sigma(\mathrm{W}_{1}) \cup \sigma(\mathrm{W}_{2})\)。（利用 b)。）
\item
  设 W 是 A′ 的一个弱闭不变向量子空间。则 \(\sigma(\mathrm{W})\) 是由这样的 \(\chi \in \widehat{A}\) 组成的集合：条件 \(x * x' = 0\)（对一切 \(x' \in W\)）蕴含 \(\chi(x) = 0\)。
\item
  若 A 有单位元，则 \(A'\) 的每个弱闭不变向量子空间都含有一个非零特征标。
\item
  在本习题的其余部分，假设 A 是无根、含单位元的正则代数。若 W 是 A′ 的一个弱闭不变向量子空间，而 U 是 \(\sigma(\mathrm{W})\) 中 \(\mathsf{X}(\mathsf{A})\) 的一个邻域，则 W 含于由 U 生成的 A′ 的弱闭向量子空间之中。
\end{enumerate}

\begin{enumerate}
\def\labelenumi{\alph{enumi})}
\setcounter{enumi}{6}
\tightlist
\item
  设 \(x' \in A'\)，并设 W 为由 \(A'\) 生成的 \(x'\) 的弱闭不变向量子空间。令 \(\sigma(x') = \sigma(\mathrm{W})\)。若 \(x \in A\) 使 \(\mathcal{G}x = 1\) 在 \(\sigma(x')\) 的一个邻域上成立，则 \(x \star x' = x'\)。（利用 \(f\)。）
\end{enumerate}

\begin{enumerate}
\def\labelenumi{\alph{enumi})}
\setcounter{enumi}{7}
\tightlist
\item
  若 \(\sigma(x')\) 是两个互不相交的闭集 \(\sigma_1\) 与 \(\sigma_2\) 的并，则存在 \(x_1', x_2' \in A'\) 使 \(\sigma(x_1') = \sigma_1\)、\(\sigma(x_2') = \sigma_2\) 且 \(x' = x_1' + x_2'\)。（令 \(x_1' = u_1 \star x'\)、\(x_{2}^{\prime}=u_{2}\star x^{\prime}\)，其中 \(\mathcal{G}u_{1}=1\) 在 \(\sigma_{1}\)（相应地，\(\sigma_{2}\)）的邻域中取 1（相应地，0），而 \(\mathcal{G}u_{2}=1\) 在 \(\sigma_{2}\)（相应地，\(\sigma_{1}\)）的邻域中取 1（相应地，0）的。）
\end{enumerate}

\begin{enumerate}
\def\labelenumi{\arabic{enumi})}
\setcounter{enumi}{3}
\item
  设 A 为满足 Ditkin 条件的交换正则半单 Banach 代数，I 为 A 的一个闭理想，\(x \in \Upsilon(\mathrm{V}(\mathrm{I}))\)，F 为 \(\mathrm{V}(Ax)\) 的边界。若 \(\mathrm{F} \cap \mathrm{V}(\mathrm{I})\) 不含非空的完满集，则 \(x \in I\)。（模仿 I，第 94 页命题 5 的推理。）
\item
  设 A 为无根的正则交换 Banach 代数。假设 A 的每个闭理想都是一些正则极大理想的交。设 \(x \in A\)。设 F 为 \(\mathcal{G}'(x)\) 的零点集。则 x 是形如 \(u_{n}x\) 的元素的极限，其中 \(u_{n} \in A\) 且 \(\mathcal{G}(u_{n})\) 在 F 的一个邻域中为零。（设 J 为使 \(y \in A\) 具有与 F 不相交的紧支撑的 \(\mathcal{G}(y)\) 所成的集合。我们有 \(x \in \overline{J}\)。利用 I，第 91 页的命题 4。）特别地，A 满足 Ditkin 条件。
\item
  设 B 为 {[}0, 1{]} 上连续可微复函数所成的代数，赋以范数 \(\|f\| = \sup |f| + \sup |f'|\)。设 A 为 B 中满足 \(f \in \mathrm{B}\) 的 \(f\left(\frac{1}{2}\right) = 0\) 所成的闭子代数。设 I 为 A 中满足 \(f \in \mathrm{A}\) 的 \(f'\left(\frac{1}{2}\right) = 0\) 所成的闭理想。
\end{enumerate}

\begin{enumerate}
\def\labelenumi{\alph{enumi})}
\item
  空间 X(A) 典范地等同于 {[}0, 1{]} = \{ \(\frac{1}{2}\) \}，且 A 是正则的。
\item
  理想 I 是 A 的一个极大理想，但不是正则的。
\end{enumerate}

\begin{enumerate}
\def\labelenumi{\arabic{enumi})}
\setcounter{enumi}{6}
\tightlist
\item
  设 R、\(T_{1}\)、\(T_{2}\) 为 \(C^{4}\) 中由满足下列条件的元素 \((z_{1}, z_{2}, z_{3}, z_{4}) \in \mathbf{C}^{4}\) 组成的子集：
\end{enumerate}

\[
\mathrm{R} \colon z _ {1} z _ {2} = 2, \qquad 1 \leqslant | z _ {1} | \leqslant 2, \qquad z _ {3} = z _ {4} = 0,
\]

\[
\mathrm{T} _ {1} \colon z _ {1} z _ {2} = 2, \qquad | z _ {1} | = 1, \qquad | z _ {3} | \leqslant 1, \qquad z _ {4} = 0,
\]

\[
\mathrm{T} _ {2} \colon z _ {1} z _ {2} = 2, \qquad | z _ {1} | = 2, \qquad | z _ {3} | \leqslant 1, \qquad z _ {4} = z _ {3} ^ {2}.
\]

\begin{enumerate}
\def\labelenumi{\alph{enumi})}
\tightlist
\item
  集合 \(X = R \cup T_{1} \cup T_{2}\) 是多项式凸的。（证明 X 是由满足下列条件的 \((z_{1}, z_{2}, z_{3}, z_{4})\) 组成的集合：
\end{enumerate}

\[
z _ {1} z _ {2} = 2, \quad | z _ {1} | \leqslant 2, \quad | z _ {2} | \leqslant 2, \quad | z _ {3} | \leqslant 1, \quad z _ {4} (z _ {4} - z _ {3} ^ {2}) = 0,
\]

\[
| z _ {4} z _ {2} ^ {k} | \leqslant 1 \quad \text { and } \quad | z _ {1} ^ {k} (z _ {4} - z _ {3} ^ {2}) | \leqslant 1 \quad \text { for } \quad k = 1, 2, \ldots).
\]

\begin{enumerate}
\def\labelenumi{\alph{enumi})}
\setcounter{enumi}{1}
\item
  设 A 为 \(\mathcal{C}(\mathrm{X})\) 的、由 \(\mathbf{C}^4\) 上多项式函数在 X 上的限制所生成的闭子代数。则 X 典范地等同于 \(\mathsf{X}(\mathsf{A})\)。
\item
  设 \(\Gamma_1, \Gamma_2 \subset \mathbf{C}\) 为以 0 为心、半径分别为 1、2 且按正向定向的圆周。设 \(\mu_0, \mu_1, \nu\) 为 \(\Gamma_1, \Gamma_2, \Gamma_1\) 上由微分形式 \(z^{-1} dz, -z^{-1} dz, z^{-2} dz\) 所定义的测度。令 \(\mu = \mu_0 + \mu_1\)，于是 \(\mu \otimes \nu\) 是 \((\Gamma_1 \cup \Gamma_2) \times \Gamma_1\) 上的一个测度。设 B 为由使 \(z = (z_1, z_2, z_3, z_4) \in X\) 满足 \((z_1, z_3) \in (\Gamma_1 \cup \Gamma_2) \times \Gamma_1\) 的 z 所组成的集合。映射 \(\varphi: z \mapsto (z_1, z_3)\)（从 B 到 \((\Gamma_{1} \cup \Gamma_{2}) \times \Gamma_{1}\) 中）是一个同胚。设 \(\nu\) 为测度 \(\varphi^{-1}(\mu \otimes \nu)\)。证明 \(\nu\) 与 A 正交。
\item
  设 g 为 X 上的函数：在 \(R \cup T_{1}\) 上为零，而在 \(z_{3}\) 上等于 \(T_{2}\)。则 \(g \notin A\)。
\item
  函数 g 在每个 \(z \in X\) 的一个邻域中与 A 的一个元素相合。
\end{enumerate}

\exercisesection{星代数}

\begin{enumerate}
\def\labelenumi{\arabic{enumi})}
\item
  找出一个幺对合 Banach 代数 A 和一个酉元 \(u \in A\)，使 \(\|u\| \neq 1\)。
\item
  设 A 为 Banach 代数，\(A_{1} \subset A\) 为其闭子代数。若 \(A_{1}\) 是对合代数，则 \(A_{1}\) 的每个特征标都可以延拓为 A 的特征标。
\item
  设 A 为幺对合代数。假设对 A 中的每个 x，x 的谱都不含 −1。
\end{enumerate}

\begin{enumerate}
\def\labelenumi{\alph{enumi})}
\item
  对 A 的每个幂等元 x，证明存在 A 的一个厄米幂等元 y，使 \(xA = yA\)。（考虑 \(1 + (x - x^{*})(x^{*} - x)\) 的逆 z；证明 z 与 x 交换，并令 \(y = xx^{*}z\)。）
\item
  对 A 的每个幂等元 \(x\)，证明存在 A 的一个厄米幂等元 \(y\) 与 A 中可逆的 \(u\)，使 \(x = uyu^{-1}\)。（若 \(x\mathrm{A} = y\mathrm{A}\)，则令 \(u = (1 - (y - x))^{-1}\)。）
\end{enumerate}

\begin{enumerate}
\def\labelenumi{\arabic{enumi})}
\setcounter{enumi}{3}
\item
  设 A 为赋范代数，带有一个连续对合。若对一切 \(\|x\|' = \sup(\|x\|, \|x^{*}\|)\) 令 \(x \in A\)，则 A 成为一个对合赋范代数，且新范数与旧范数等价。
\item
  设 A 为无根交换 Banach 代数。A 的每个对合都是连续的。（利用 I，第 40 页的命题 9。）
\end{enumerate}

\begin{enumerate}
\def\labelenumi{¶ \arabic{enumi})}
\setcounter{enumi}{5}
\tightlist
\item
  设 A 为每个特征标都是厄米的交换幺对合 Banach 代数。假设对每个不可逆的厄米元 \(x \in A\)，有 \(\|(x - \lambda)^{-1}\| = o(\mathcal{I}(\lambda)^{-2})\)（当 \(\lambda\) 保持 \(\mathcal{I}(\lambda) \neq 0\) 而趋于 0 时）。则 A 的每个闭理想 I 都是一些极大理想的交。（设 \(\varphi: A \longrightarrow A/I\) 为典范态射。设 \(x \in A\) 使 \(\varphi(x)\) 属于 A/I 的根。要证明 \(\varphi(x) = 0\)。化归到 \(x\) 是厄米的情形。然后利用 I，第 161 页的习题 29。）
\end{enumerate}

\begin{enumerate}
\def\labelenumi{\arabic{enumi})}
\setcounter{enumi}{6}
\item
  设 A 为交换幺对合 Banach 代数。为使 A 的每个特征标都是厄米的，充要条件是对每个 \(1 + xx^{*}\)，\(x \in A\) 可逆。（若 A 的每个特征标都厄米，则 \(\mathcal{G}(1 + xx^{*})\) 在 \textgreater{} 0（在 \(\mathsf{X}(\mathsf{A})\) 上）处处成立。若 A 的特征标 \(\chi\) 非厄米，则存在厄米元 \(x \in A\) 使 \(\chi(x) = i\)，从而 \(\chi(1 + xx^{*}) = 0\)，而 \(1 + xx^{*}\) 不可逆。）
\item
  构造一个群 G 的例子，使对合 Banach 代数 \(\mathcal{M}^{1}(G)\)（I，第 100 页，例 4）不是星代数。
\item
  设 A 为幺 Banach 代数。A 的酉元所成的群是 A* 的一个极大有界子群。
\item
  设 A 为对合 Banach 代数，\(\widetilde{\mathsf{A}}\) 为添加单位元后所得的代数。证明：I，第 15 页上在 \(\widetilde{\mathsf{A}}\) 上定义的范数，一般说来与 I，第 105 页命题 3 所定义的范数不相一致。
\item
  设 A 为对合 Banach 代数，其对合满足：对一切 \(\|x^{*}x\| = \|x\| \cdot \|x^{*}\|\) 有 \(x \in A\)。则 A 是一个星代数。（对厄米元 y 有 \(\|y^{2}\| = \|y\|^{2}\)，从而 \(\|y\| = \varrho(y)\)；注意到对一切 \(\mathrm{Sp}_{\mathrm{A}}^{\prime}(x^{*}) = \overline{\mathrm{Sp}_{\mathrm{A}}^{\prime}(x)}\) 有 \(x \in A\)，因此我们有
\end{enumerate}

\[
\| x ^ {*} \| \cdot \| x \| = \| x ^ {*} x \| = \varrho (x ^ {*} x) \leqslant \varrho (x ^ {*}) \varrho (x) = \varrho (x ^ {2}) \leqslant \| x \| ^ {2},
\]

从而 \(\|x^{*}\|\leqslant\|x\|\) 且 \(\|x^{*}\|=\|x\|\)。

\begin{enumerate}
\def\labelenumi{\arabic{enumi})}
\setcounter{enumi}{11}
\item
  设 A 为星代数，N 为 A 的正规元所成的集合，\(x \in N\)，V 为 C 中 0 的邻域。存在 x 在 N 中的一个邻域 U，使得对每个 \(y \in U\)，有 \(\mathrm{Sp}_{\mathrm{A}}(y) \subset \mathrm{Sp}_{\mathrm{A}}(x) + \mathrm{V}\) 与 \(\mathrm{Sp}_{\mathrm{A}}(x) \subset \mathrm{Sp}_{\mathrm{A}}(y) + \mathrm{V}\)。（利用 I，第 108 页命题 1 的推论、I，第 23 页的命题 3 以及 I，第 76 页的命题 10。）
\item
  设 A 为 C 上的幺 Banach 代数，\(x \in A\)，且 \(S = \mathrm{Sp}_{\mathrm{A}}(x)\)。假设存在从 \(\varphi\) 到 A 中的一个态射 \(\mathcal{C}(\mathrm{S})\)，它把 1 映为 1、把 z 映为 x（z 表示 S 上的恒等映射）。则对每个 \(\|\varphi(f)\| \geqslant \|f\|\)，有 \(f \in \mathcal{C}(\mathrm{S})\)。特别地，若态射 \(\varphi\) 连续，则 \(\varphi\) 是到其像上的一个同胚。（可以假设 A 交换。设 \(\lambda \in S\)。存在 \(\chi \in \mathrm{X}(\mathrm{A})\) 使 \(\lambda = \chi(x) = \chi(\varphi(z)) = (\mathrm{X}(\varphi)(\chi))(z)\)。于是 \(\mathrm{X}(\varphi)(\chi)\) 就是 \(f \mapsto f(\lambda)\) 的特征标 \(\mathcal{C}(\mathrm{S})\)。设 \(f \in \mathcal{C}(\mathrm{S})\)。由前所证，存在 \(\chi \in \mathrm{X}(\mathrm{A})\) 使 \(|f|\) 在 \(\mathrm{X}(\varphi)(\chi)\) 处达到其最大值。于是 \(|\chi(\varphi(f))| = \|f\|\)，从而 \(\|\varphi(f)\| \geqslant \|f\|\)。）
\item
  设 G 为局部紧群。若 Stell(G) 有单位元，则群 G 是离散的。（利用下述事实：在 Hilbert 空间 \(L^{2}(G)\) 中，恒等映射是诸自同态 \(\gamma_{f}\) 的范数极限，其中 \(\gamma\) 是左正则表示，而 \(f \in L^{1}(G)\)。）
\item
  \begin{enumerate}
  \def\labelenumii{\alph{enumii})}
  \tightlist
  \item
    设 A 为交换星代数。若 x 与 y 是 A 中满足 \(x \leqslant y\) 的正元素，则 \(x^{2} \leqslant y^{2}\)。
  \end{enumerate}
\end{enumerate}

\begin{enumerate}
\def\labelenumi{\alph{enumi})}
\setcounter{enumi}{1}
\item
  存在一个星代数 A 及 A 中元素 x、y，使得 \(0 \leqslant x \leqslant y\)，但 \(x^{2}\) 并不成立\(\leqslant y^{2}\)。
\item
  存在一个星代数 A 及 A 中元素 \(x\) 与 \(y\)，使得 \(|x + y|\) 并不小于等于 \(\leqslant |x| + |y|\)。
\end{enumerate}

\begin{enumerate}
\def\labelenumi{\arabic{enumi})}
\setcounter{enumi}{15}
\tightlist
\item
  设 X 为拓扑空间。设 \(\mathcal{C}_{b}(\mathrm{X})\) 为 X 上有界连续复值函数所成的交换星代数。
\end{enumerate}

\begin{enumerate}
\def\labelenumi{\alph{enumi})}
\item
  设 \(\widehat{\mathrm{X}} = \mathrm{X}(\mathcal{C}_b(\mathrm{X}))\)；空间 \(\widehat{\mathrm{X}}\) 是一个紧空间，使得 \(\mathcal{C}(\widehat{\mathrm{X}})\) 典范地等同于 \(\mathcal{C}_b(\mathrm{X})\)。
\item
  映射 \(j: X \to \widehat{X}\)（把 x 对应到特征标 \(f \mapsto f(x)\)）是连续的。
\item
  若 X 是完全正则的（TG，IX，第 8 页，定义 4），则 j 是 X 到 \(\widehat{X}\) 的一个稠密子空间上的同胚。（为证 j 的像在 \(\widehat{X}\) 中稠密，证明：在 \(\widehat{X}\) 上连续且在 \(j(\mathrm{X})\) 上为零的复值函数恒为零。）
\item
  空间 \(\widehat{X}\) 典范地等同于 X 的 Stone–Čech 紧化（TG，IX，第 9 页）。
\end{enumerate}

\begin{enumerate}
\def\labelenumi{\arabic{enumi})}
\setcounter{enumi}{16}
\item
  设 A 为 C*-代数，\(x \in \mathrm{A}\)。若 \(x = a - b\)（其中 \(a, b\) 为正元素且 \(ab = ba = 0\)），则 \(a = x^{+}\) 且 \(b = x^{-}\)。
\item
  设 A 为可分星代数。存在 \((e_{n})_{n\in\mathbf{N}}\) 中的一个序列 \(A_{+}\)，使相联的基本滤子（TG，I，第 42 页，定义 7）是 A 的一个递增近似单位元。
\item
  设 A 为星代数。称元素 \(x\) 是系统正的，如果 \(x \in \mathrm{A}_{+}\) 且 \(x\mathrm{A}x\) 在 A 中稠密。
\end{enumerate}

\begin{enumerate}
\def\labelenumi{\alph{enumi})}
\item
  证明：交换星代数 A 含有系统正元素，充要条件是存在一个在无穷远处可数的局部紧拓扑空间 X，使 A 同构于 \(\mathcal{C}_{0}(X)\)。
\item
  若 A 是幺的，则 A 的元素 x 是严格正的，当且仅当 x 是正的且可逆。
\item
  设 \(x \in A\) 是一个系统正元素，满足 \(\|x\| \leqslant 1\)。证明序列 \((x^{1/n})_{n \geqslant 1}\) 是 A 中的一个递增近似单位元。
\item
  星代数 A 含有系统正元素，充要条件是存在 A 中的一个序列 \((x_{n})_{n\geqslant1}\)，它构成 A 的一个近似恒等元。
\end{enumerate}

\begin{enumerate}
\def\labelenumi{\arabic{enumi})}
\setcounter{enumi}{19}
\item
  设 A 为 C*-代数，I 为 A 的一个闭双边理想，B 为 A 的一个闭对合子代数。则 B+I 是 A 的一个闭对合子代数，且存在等距同构 \(B + I/I \rightarrow B/B \cap I\)。
\item
  设 S 为 C 中的闭单位圆盘。设 A 为从 S 到 C 的、在开单位圆盘中全纯的连续函数所成的对合 Banach 代数，其对合由 \(f^{*}(z) = \overline{f(\overline{z})}\) 给出。A 中存在一个厄米元素，其谱不含于 R 中。
\item
  设 \(\mathrm{A} = \mathcal{C}([0,1])\)。证明由 {[}0, 1{]} 上的恒等函数所生成的 A 的理想 I 不是闭的。
\item
  设 A 为幺 C*-代数。用 \(A_{+}^{\prime}\) 表示 A 上的线性形式 f（不必连续）所成的向量空间，使得
\end{enumerate}

\[
f (\mathrm{A} _ {+}) \subset [ 0, + \infty [
\]

（“正线性形式”）。

\begin{enumerate}
\def\labelenumi{\alph{enumi})}
\item
  设 \(f \in A_{+}^{\prime}\)。对 \((x_{n})_{n \in \mathbf{N}}\) 中的每个序列 \(A_{+}^{\leqslant 1}\)，族 \((f(x_{n}))\) 有界。（证明：对 \(\sum_{i} f(x_{i})y_{i}\) 中的每个正族 \((y_{i})\)，级数 \(\ell^{1}(\mathbf{N})\) 收敛。）
\item
  线性映射 f 是连续的。
\item
  线性形式 \(f\) 属于 \(\mathrm{A}_{+}^{\prime}\) 的充要条件是 \(f\) 连续且 \(\| f \| = f(1)\)。
\item
  设 \(x \in A_{h}\) 为厄米元素。存在正线性形式 \(f \in A_{+}^{\prime}\)，使 \(|f(x)| = \|x\|\)。（利用 I，第 174 页的习题 7。）
\end{enumerate}

\begin{enumerate}
\def\labelenumi{\arabic{enumi})}
\setcounter{enumi}{23}
\tightlist
\item
  设 A 为 C 上的幺 Banach 代数。假设 \(x \mapsto x^{*}\) 与 \(x \mapsto x^{\circ}\) 是 A 上的对合，且 A 赋以其中任一对合都成为星代数。用 \(A_{1}\) 与 \(A_{2}\) 记这两个星代数。
\end{enumerate}

\begin{enumerate}
\def\labelenumi{\alph{enumi})}
\item
  沿用习题 23 的记号，证明 \(\mathrm{A}_{1,+}^{\prime} = \mathrm{A}_{2,+}^{\prime}\)。用 \(\mathrm{A}_{+}^{\prime}\) 记这个集合。
\item
  设 \(x \in A_{1,h}\) 是满足 \(x^{*} = x\) 的元素。证明：对每个 \(f \in A_{+}^{\prime}\)，有 \(f(i(x - x^{\circ})) = 0\)。由此推出 \(x \in A_{2,h}\)。（利用习题 23 的 b)。）
\item
  由此得出结论：\(A_{1} = A_{2}\)，并且一个复 Banach 代数至多有一个使它成为 C*-代数的对合。
\end{enumerate}

\begin{enumerate}
\def\labelenumi{\arabic{enumi})}
\setcounter{enumi}{24}
\tightlist
\item
  设 G 为群。用 C{[}G{]} 表示 G 在 C 上的群代数。它是 Hilbert 空间 \(\ell^{2}(\mathrm{G})\) 的一个子空间。对 G 中的 g，用 {[}g{]} 表示 C{[}G{]} 中对应于 g 的元素。诸 {[}g{]}（对 \(g \in G\)）构成 \(\ell^{2}(G)\) 的一个标准正交基，且 C{[}G{]} 在 \(\ell^{2}(G)\) 中稠密。
\end{enumerate}

\begin{enumerate}
\def\labelenumi{\alph{enumi})}
\tightlist
\item
  映射
\end{enumerate}

\[
\sum_ {g} \alpha_ {g} [ g ] \mapsto \sum_ {g} \overline {{\alpha}} _ {g} [ g ^ {- 1} ]
\]

是 \(\mathbf{C}[\mathrm{G}]\) 上的一个对合。

\begin{enumerate}
\def\labelenumi{\alph{enumi})}
\setcounter{enumi}{1}
\tightlist
\item
  对 \(g \in G\)，设 \(\pi(g)\) 为 C\(x \mapsto gx\)[{G}]{ 的映射 }。把 g 对应到 \(\pi(g)\) 的映射可延拓为从 C{[}G{]} 到 \(\mathcal{L}(\ell^{2}(G))\) 中的一个连续单射对合代数态射。
\end{enumerate}

用 \(\operatorname{Stell}_{r}(\mathbf{G})\) 记在 \(\mathcal{L}(\ell^{2}(\mathbf{G}))\) 中映射 \(\pi\) 的像的闭包，称之为 G 的约化 C*-代数。它是一个 C*-代数。我们把 \(\mathbf{C}[\mathbf{G}]\) 经 \(\pi\) 在 \(\operatorname{Stell}_{r}(\mathbf{G})\) 中的像与它等同。

\begin{enumerate}
\def\labelenumi{\alph{enumi})}
\setcounter{enumi}{2}
\tightlist
\item
  线性迹映射 \(\operatorname{Tr} \colon \mathbf{C}[\mathrm{G}] \to \mathbf{C}\)（满足 \(\operatorname{Tr}(1) = 1\) 且对 \(\operatorname{Tr}(g) = 0\) 有 \(g \neq 1\)）可延拓为连续线性映射 \(\operatorname{Stell}_r(\mathrm{G}) \to \mathbf{C}\)。
\end{enumerate}

\begin{enumerate}
\def\labelenumi{\alph{enumi})}
\setcounter{enumi}{3}
\tightlist
\item
  我们有 \(\operatorname{Tr}(xy) = \operatorname{Tr}(yx)\)（对一切 \(x\) 与 \(y\)，它们都属于 \(\operatorname{Stell}_r(G) \)）。
\end{enumerate}

\begin{enumerate}
\def\labelenumi{\alph{enumi})}
\setcounter{enumi}{4}
\item
  我们有：对一切 \(\operatorname{Tr}(xx^{*}) \geqslant 0\) 属于 \(x\)，\(\operatorname{Stell}_r(G)\)；且 \(\operatorname{Tr}(xx^{*}) = 0\) 当且仅当 \(x = 0\)。
\item
  设 e 是 C{[}G{]} 的一个幂等元，满足 \(e \neq 0\) 且 \(e \neq 1\)。证明 \(0 < \operatorname{Tr}(e) < 1\)。（利用习题 3。）
\item
  设 e 是 Z{[}G{]} 的一个幂等元。证明 e = 0 或 e = 1。（“Kaplansky 定理”。）
\item
  设 x 是 C{[}G{]} 的一个元素。则 x 可逆，当且仅当它左可逆，当且仅当它右可逆。
\item
  设 n 为严格正的整数。类型为 \((n, n)\) 且系数在 C{[}G{]} 中的矩阵 x 可逆，当且仅当它左可逆，当且仅当它右可逆。
\end{enumerate}

\begin{enumerate}
\def\labelenumi{\arabic{enumi})}
\setcounter{enumi}{25}
\item
  设 A 为幺 C*-代数。设 \(a\) 与 \(b\) 为 A 的正规元，\(x \in \mathrm{A}^*\) 满足 \(b = xax^{-1}\)。存在 A 的酉元 \(u\) 使 \(b = uau^*\)。（利用前一习题以及 \(x\) 的极分解，注意 \(|x|\) 属于 \(x\) 的双交换子之中。）
\item
  设 A 为幺 C*-代数，a、b、c 为 A 的元素。假设 b 与 c 正规。若 ac = cb，则对 a 的谱与 b 的谱的并上的每个连续函数 f，有 \(f(a)c = cf(b)\)。
\item
  设 A 为 C*-代数。设 \(a\) 与 \(b\) 为 A 的正元素，满足 \(ab \neq 0\)。则 \(\mathrm{Sp}_{\mathrm{A}}(ab)\) 不化为 0。
\item
  \begin{enumerate}
  \def\labelenumii{\alph{enumii})}
  \tightlist
  \item
  设 G 为局部紧拓扑群。Banach 代数 L\(^{1}\)(G) 具有一个近似恒等元。（利用 INT，VIII，§2，第 7 节，引理 4。）
  \end{enumerate}
\end{enumerate}

\begin{enumerate}
\def\labelenumi{\alph{enumi})}
\setcounter{enumi}{1}
\item
  设 A 为具有近似恒等元的 Banach 代数。对每个特征标 \(\chi\in\mathsf{X}(\mathrm{A})\)，有 \(\|\chi\|=1\)。
\item
  L\(^{1}\)(Z) 的由序列（\(x_{n}\)）——满足 \(x_{n} = 0\)（当 \(n < 0\)）——所组成的 Banach 子代数没有近似恒等元。（构造一个范数 \(< 1\) 的非零特征标。）
\end{enumerate}

\begin{enumerate}
\def\labelenumi{\alph{enumi})}
\setcounter{enumi}{3}
\tightlist
\item
  设 G 为由两个元素 \(a\) 与 \(b\) 生成的自由群，\(\mu\) 为 G 上的一个 Haar 测度。设 I 为 \(\mathrm{L}^1 (\mathrm{G})\) 中由使 \(f\in \mathrm{L}^{1}(\mathrm{G})\) 的 \(\mu (f) = 0\) 所成的闭理想。令 \(\xi = e^{2i\pi /3}\)。令 \(x = \varepsilon_{e} + \xi \varepsilon_{a} + \xi^{-1}\varepsilon_{b}\in \mathrm{L}^{1}(\mathrm{G})\)。对每个 \(f\in \mathrm{L}^{1}(\mathrm{G})\)，有 \(\| x*f - x\| \geqslant 1\)。Banach 代数 I 没有近似恒等元。
\end{enumerate}

\begin{enumerate}
\def\labelenumi{¶ \arabic{enumi})}
\setcounter{enumi}{29}
\tightlist
\item
  设 A 为具有近似恒等元的 Banach 代数。用 \(\widetilde{\mathsf{A}}\) 表示由 A 添加单位元（记作 1）所得的 Banach 代数。
\end{enumerate}

\begin{enumerate}
\def\labelenumi{\alph{enumi})}
\item
  设 \(x \in \mathrm{A}\)。则 \(x\) 属于由 \(x\) 生成的闭左理想。
\item
  对每个 \(\varepsilon > 0\) 与 \((x_i)_{i \in \mathrm{I}}\) 中元素的每个有限族 \(\mathrm{A}\)，存在 \(e \in \mathrm{A}\) 使 \(\| e \| \leqslant 1\)，且对一切 \(\| ex_i - x_i \| < \varepsilon\) 有 \(i \in \mathrm{I}\)。
\item
  设 \(x \in A\)、\(\varepsilon > 0\) 与 \(0 < \gamma \leqslant 1/8\)。存在 \(\eta > 0\)，使得对满足 \(e \in A\) 与 \(\|e\| \leqslant 1\) 的每个 \(\|ex - x\| < \eta\)，元素 \(\alpha = (1 - \gamma) \cdot 1 + \gamma e\) 在 \(\widetilde{A}\) 中可逆，且 \(\|x - \alpha^{-1}x\| < \varepsilon\)。
\end{enumerate}

\begin{enumerate}
\def\labelenumi{\alph{enumi})}
\setcounter{enumi}{3}
\tightlist
\item
  设 \(x \in \mathrm{A}\)、\(\varepsilon > 0\)、\(\delta > 0\) 与 \(0 < \gamma \leqslant 1/8\)。存在 A 中满足下列性质的序列 \((e_n)_{n \geqslant 1}\)：

\begin{enumerate}
\def\labelenumi{(\roman{enumi})}
\item
  对一切 \(\|e_{n}\| < 2\)，有 \(n \geqslant 1\)；
\item
  元素
\end{enumerate}

\[
y _ {n} = \sum_ {k = 1} ^ {n} \gamma (1 - \gamma) ^ {k - 1} e _ {k} + (1 - \gamma) ^ {n} \cdot 1
\]

在 \(\widetilde{\mathrm{A}}\) 中（对一切 \(n\geqslant 1\)）可逆，且

\[
\left\| y _ {n} ^ {- 1} x - y _ {n - 1} ^ {- 1} x \right\| <   \delta 2 ^ {- n} \quad (n \geqslant 2), \quad \left\| y _ {1} ^ {- 1} x - x \right\| <   \delta / 2.
\]

（假设 \((e_{1},\ldots,e_{n})\) 已经定义，考虑元素 \(e_{n+1}\)（任意，范数 \textless{} 2）；注意

\[
y _ {n + 1} = ((1 - \gamma) \cdot 1 + e _ {n + 1}) \left(\sum_ {k = 1} ^ {n} \gamma (1 - \gamma) ^ {k - 1} e _ {k} ^ {\prime} + (1 - \gamma) ^ {n} \cdot 1\right),
\]

其中 \(e_k' = ((1 - \gamma) \cdot 1 + e_{n+1})^{-1} e_k\)，并确定满足所需条件的 \(e_{n+1}\)（利用 \(\widetilde{\mathsf{A}}\) 的可逆元所成集合是开集、且映射 \(a \mapsto a^{-1}\) 连续这些事实）。
\end{enumerate}

\begin{enumerate}
\def\labelenumi{\alph{enumi})}
\setcounter{enumi}{4}
\tightlist
\item
  对每个 \(x \in A\)，存在 A 中的 y 与 z 使 x = yz（“Cohen 因式分解定理”）。设 \(\delta > 0\)。可以假设 z 属于由 \(x\) 生成的闭左理想，且 \(\|y-z\|<\delta\)。（考虑如上的序列 \((e_{n})_{n\geqslant1}\)，并按之前的方式定义 \(y_{n}\)；序列 \(z_{n}=y_{n}^{-1}x\) 收敛于一个元素 \(z\in\mathrm{A}\)，且 \(z_{n}\) 属于由 \(x\) 生成的闭左理想；序列 \((y_{n})\) 收敛于 \(y\in\mathrm{A}\)，且 \(x=yz\)。）
\end{enumerate}

\begin{enumerate}
\def\labelenumi{\arabic{enumi})}
\setcounter{enumi}{30}
\tightlist
\item
  设 A 为关于 Lebesgue 测度在 {[}0, 1{]} 上可积的复函数所成的 Banach 空间。令
\end{enumerate}

\[
(f \cdot g) (t) = \int_ {0} ^ {t} f (t - x) g (x) d x
\]

对 \(t \in [0, 1]\)。

\begin{enumerate}
\def\labelenumi{\alph{enumi})}
\item
  函数 \(f \cdot g\) 几乎处处有定义，且是 A 的一个元素。
\item
  赋以运算 \((f,g)\mapsto f\cdot g\) 后，空间 A 是一个交换 Banach 代数。
\item
  设 \(x \in \mathrm{A}\) 是恒等于 1 的常值函数。函数 \(x^n\) 是
\end{enumerate}

\[
t \mapsto \frac {t ^ {n - 1}}{(n - 1) !}.
\]

\begin{enumerate}
\def\labelenumi{\alph{enumi})}
\setcounter{enumi}{3}
\item
  代数 A 由 x 拓扑地生成，元素 x 是拟幂零元，且 A 等于它的根。
\item
  代数 A 具有近似恒等元。
\item
  代数 A 没有极大理想，无论是正则的还是非正则的，无论闭与否（利用前一习题以及 I，第 166 页的习题 1）。特别地，X(A) 是空集，因而是紧的。
\end{enumerate}

\begin{enumerate}
\def\labelenumi{\arabic{enumi})}
\setcounter{enumi}{31}
\tightlist
\item
  设 G 为群。称 G 具有 Powers 性质，如果对每个 \(n \in N\) 以及 G = \(G = \{e\}\) 的每个有限子集 F，存在 G 分成两个子集 A 与 B 的一个划分，以及 G 的元素的一个族 \((g_i)_{1 \leqslant i \leqslant n}\)，使得
\end{enumerate}

\begin{enumerate}
\def\labelenumi{(\roman{enumi})}
\item
  对一切 \(g \in \mathrm{F}\)，集合 A 与 gA 不相交；
\item
  对任意两个不同的整数 \(j\) 与 \(k\)（介于 1 与 \(n\) 之间），集合 \(g_{j}\mathrm{B}\) 与 \(g_{k}\mathrm{B}\) 不相交。
\end{enumerate}

设 G 为具有 Powers 性质的群。设 I 为约化 C*-代数 \(\text{Stell}_{r}(\text{G})\)（习题 25）的一个非零闭双边理想。设 \(\pi: C[G] \to \text{Stell}_{r}(G)\) 为典范表示。

\begin{enumerate}
\def\labelenumi{\alph{enumi})}
\tightlist
\item
  存在一个族 \((\lambda_g)_{g\in \mathrm{G - }\{e\}}\)，其值属于 \(\mathbf{C}\)，使得
\end{enumerate}

\[
x = \sum_ {g \in \mathrm{G} - \{e \}} \lambda_ {g} \pi (g)
\]

是 \(\operatorname{Stell}_{r}(\mathrm{G})\) 的一个厄米元素，且 \(y = 1 + x\) 属于 I。

\begin{enumerate}
\def\labelenumi{\alph{enumi})}
\setcounter{enumi}{1}
\tightlist
\item
  设 \(\varepsilon > 0\) 满足 \(\varepsilon < 1\)。存在 G = \{e\} 的一个有限子集 F 以及一个整数 \(n \geqslant 1\)，使得元素
\end{enumerate}

\[
x ^ {\prime} = \sum_ {g \in \mathrm{F}} \lambda_ {g} \pi (g)
\]

满足 \(\|x' - x\| < \varepsilon\) 与 \(\|x'\| < \frac{\sqrt{n}}{2}(1 - \varepsilon)\)。

\begin{enumerate}
\def\labelenumi{\alph{enumi})}
\setcounter{enumi}{2}
\tightlist
\item
  取 \(G = A \cup B\) 与 \((g_{i})_{1 \leqslant i \leqslant n}\)，它们对有限子集 F 与整数 n 满足 Powers 性质的条件。令
\end{enumerate}

\[
y ^ {\prime} = \frac {1}{n} \sum_ {i = 1} ^ {n} \pi (g _ {i}) x ^ {\prime} \pi (g _ {i} ^ {- 1}).
\]

我们有 \(\|y'\| < 1 - \varepsilon\)。（设 \(p_i\) 是 \(\ell^{2}(\mathrm{G})\) 到 \(\ell^{2}(g_{i}\mathrm{B})\) 上的正交投影；验证

\[
(1 - p _ {i}) \pi (g _ {i}) x ^ {\prime} \pi (g _ {i} ^ {- 1}) (1 - p _ {i}) = 0
\]

对 \(1 \leqslant i \leqslant n\) 成立，并由此推出

\[
1 + y ^ {\prime} = 1 + \frac {1}{n} \sum_ {i = 1} ^ {n} p _ {i} \pi (g _ {i}) x ^ {\prime} \pi (g _ {i} ^ {- 1}) + \left(\frac {1}{n} \sum_ {i = 1} ^ {n} p _ {i} \pi (g _ {i}) x ^ {\prime} \pi (g _ {i} ^ {- 1}) (1 - p _ {i})\right) ^ {*},
\]

然后注意自同态 \(p_{i}\pi(g_{i})x^{\prime}\pi(g_{i}^{-1})\) 具有两两正交的像，以此估计它们的范数。

\begin{enumerate}
\def\labelenumi{\alph{enumi})}
\setcounter{enumi}{3}
\tightlist
\item
  理想 I 包含元素
\end{enumerate}

\[
1 + \frac {1}{n} \sum_ {i = 1} ^ {n} \pi (g _ {i}) x \pi (g _ {i} ^ {- 1}),
\]

并且这个元素可逆。由此推出，星代数 \(\text{Stell}_{r}(\text{G})\) 除 \(\{0\}\) 与 \(\text{Stell}_{r}(\text{G})\) 外不含任何闭双边理想。

\begin{enumerate}
\def\labelenumi{\alph{enumi})}
\setcounter{enumi}{4}
\tightlist
\item
  设 G 为非交换自由群。则 G 具有 Powers 性质。
\end{enumerate}
